\documentclass[12pt,a4paper]{article}

\usepackage[utf8]{inputenc}
\usepackage[english]{babel}
\usepackage[english]{tdk}        % formal requirements of the TDK call

\usepackage{booktabs}
\usepackage{array}
\usepackage{tabularx}
\usepackage{multirow}
\usepackage{graphicx}
\usepackage{xcolor}
\usepackage{microtype}
\usepackage{listings}
\usepackage{float}
\usepackage{enumitem}
\usepackage{pifont}
\usepackage{url}
\usepackage[hidelinks]{hyperref}
\usepackage[round,authoryear]{natbib}

\setlength{\parskip}{0.4em}
\setlength{\parindent}{0pt}

\lstset{
  basicstyle=\ttfamily\footnotesize,
  keywordstyle=\color{blue!70!black}\bfseries,
  commentstyle=\color{green!50!black}\itshape,
  stringstyle=\color{red!70!black},
  numbers=left, numberstyle=\tiny\color{gray}, numbersep=6pt,
  frame=single, framerule=0.4pt, xleftmargin=20pt, xrightmargin=5pt,
  breaklines=true, showstringspaces=false, language=Python, tabsize=4,
  captionpos=b,
}

\theoremstyle{plain}
\newtheorem{theorem}{Theorem}[section]
\newtheorem{proposition}[theorem]{Proposition}
\theoremstyle{definition}
\newtheorem{definition}[theorem]{Definition}
\newtheorem{assumption}[theorem]{Assumption}
\theoremstyle{remark}
\newtheorem{remark}[theorem]{Remark}

% ---------------------------------------------------------------- cover
\tdkev{2026}
\tdkszerzo{Marcell Moln\'ar}
\tdkszak{[PROGRAMME]}
\tdkevfolyam{[YEAR]}
\tdktemavezeto{[SUPERVISOR NAME]}
\tdkszekcio{[SECTION --- suggested: Operations Research / Transport Modelling]}
\tdkcim{A Multi-Modal Markov-Chain Passenger-Flow Model for the BKK Budapest
GTFS Network: Entropy Maximisation, Kolmogorov Equations, and Stochastic
Simulation}

\begin{document}

\tdkfrontmatter
\tdkcover

\begin{tdkabsztrakt}
\textbf{Research question.} Which stops hold the Budapest public transport
network together, and can this be answered from the published timetable alone,
without fare-collection data? The paper asks how far a first-principles
stochastic model built only from the BKK GTFS feed can describe passenger
redistribution and identify critical stops.

\textbf{Method.} The timetable is turned into three stacked stochastic layers:
a discrete-time Markov chain, a continuous-time chain whose generator combines
Wilson's entropy-maximising destination kernel with scheduled departure
frequencies, and exact Doob--Gillespie realisations with a $\tau$-leaping
approximation. Initial stop populations come from a prior anchored to published
system ridership. Criticality is measured by node removal, combining stationary
passenger throughput with the increase of Kemeny's constant and the loss of
network efficiency, using only the components that are well defined for each
chain. The same failure analysis is then inverted into a link-augmentation
problem: new metro and tram links are chosen greedily to minimise the average
and worst-case demand-weighted accessibility loss under single-station
failures, using an exact shortest-path update.

\textbf{Own results.} On 5\,177 stops the mean field agrees with exact
simulation within 1\% (metro) and with $\tau$-leaping within 2.4\% (full
network). The metro chain is periodic, so spectral-gap measures are
undefined. Criticality is concentrated in De\'ak Ferenc t\'er, K\'alvin t\'er
and Keleti p\'alyaudvar; removing De\'ak alone cuts metro efficiency by 49\%.
Three short metro links that bypass De\'ak Ferenc t\'er and K\'alvin t\'er
(R\'ak\'oczi t\'er--Arany J\'anos utca, F\H{o}v\'am t\'er--Opera,
F\H{o}v\'am t\'er--V\"or\"osmarty t\'er) cut the worst-case loss from 44\% to
19\%, about twice the robustness gain of the planned M5 inner section, while
pure line extensions add no redundancy at all.

\medskip
\noindent\emph{Keywords:} stochastic transit assignment, entropy maximisation,
continuous-time Markov chain, Gillespie algorithm, GTFS, network resilience,
Kemeny constant, network design
\end{tdkabsztrakt}


\clearpage
\tableofcontents

\tdkmaintext
% =============================================================================
\section{Introduction}
\label{sec:intro}

\subsection{Motivation and Problem Setting}

The Budapest metropolitan public transit network operated by Budapesti
K\"{o}zleked\'{e}si K\"{o}zpont (BKK) is one of Europe's denser urban transit
systems, serving on the order of $4\times10^6$ passenger journeys per weekday
across six modal families (metro, bus, tram, trolleybus, suburban railway, and
ferry).  Despite this scale, spatially resolved passenger flow---the number of
passengers at each stop at each point in time---is not directly observable from
the publicly available schedule data alone.

The practical stakes are high.  Spatial passenger density determines platform
crowding risk, vehicle overcrowding, intermodal transfer delays, and the
propagation of service disruptions.  A principled model for passenger flow
dynamics, even an imperfect one grounded in well-understood probabilistic
machinery, enables:
\begin{enumerate}
  \item \textbf{Network resilience assessment:} identifying stops and links
    whose failure disproportionately degrades service.
  \item \textbf{Prior initialisation:} providing starting distributions for
    data-assimilation procedures once partial measurements (APC counters, AFC
    tap-in records) become available.
  \item \textbf{Scenario analysis:} evaluating the system-level impact of new
    routes, closures, or timetable changes without requiring demand surveys.
\end{enumerate}

\subsection{Related Work}

\textbf{Entropy-based transport modelling.}
Wilson's~\citep{wilson1969} doubly-constrained gravity model remains the
theoretical backbone of spatial interaction modelling.  Its maximum-entropy
derivation has been formalised and extended by Erlander and
Stewart~\citep{erlander1990} and connected to random utility models by
Anas~\citep{anas1983}.  Recent work has embedded entropy transport models within
the Wasserstein optimal-transport framework~\citep{peyre2019}.

\textbf{Stochastic transit assignment.}
The logit-based stochastic user equilibrium (SUE) model of Dial~\citep{dial1971}
and the frequency-based assignment of Spiess and Florian~\citep{spiess1989} are
the industry standards.  Dynamic stochastic assignment extending to
time-varying supply is treated in Yildirimoglu and Kim~\citep{yildirimoglu2018}
and the survey of Bliemer et al.~\citep{bliemer2003}.  None of these papers
couple the assignment to a CTMC generator or uses Gillespie simulation.

\textbf{CTMC models for transit.}
Markov-chain models have been applied to pedestrian flow
\citep{helbing2000},
station waiting queues~\citep{jaiswal1961}, and subway dwell
times~\citep{zhang2019}, but full-network CTMC formulations over
GTFS-structured schedules are uncommon.  The closest precedent is the
queueing-network approach of Heidergott et al.~\citep{heidergott2010}.

\textbf{Gillespie simulation beyond biochemistry.}
Although Gillespie's SSA~\citep{gillespie1977} was developed for chemical
reaction networks, it applies to any CTMC.  Applications in
ecology~\citep{wilkinson2018}, epidemiology~\citep{allen2008}, and
traffic~\citep{maerivoet2005} demonstrate its scope.

\textbf{Network resilience.}
The node-removal framework we adopt draws on Latora and
Marchiori~\citep{latora2001}, Kemeny and Snell~\citep{kemeny1976}, and the
transit-specific vulnerability analysis of Cats and
Jenelius~\citep{cats2014}.

\subsection{Contributions and Scope}

\begin{description}[leftmargin=2em]
  \item[C1.] \textbf{CTMC generator from GTFS.}  We show precisely how the
    Wilson entropy solution translates into a CTMC generator whose units are
    dimensionally consistent and whose entries are computable from
    \texttt{stop\_times.txt}.
  \item[C2.] \textbf{Regime analysis.}  We compare the KFE mean-field
    trajectory with exact SSA and $\tau$-leap ensembles on the real network and
    characterise when each layer is the appropriate model.
  \item[C3.] \textbf{Sparse-graph formulation.}  The full $n\times n$ cost
    matrix is never formed; all operations exploit the $\ll1\%$ connectivity
    of the real GTFS graph.
  \item[C4.] \textbf{Hierarchical demand prior.}  A calibrated Poisson
    log-linear prior provides plausible initial passenger counts from freely
    available covariates.
  \item[C5.] \textbf{Resilience ranking.}  Three complementary vulnerability
    metrics (Kemeny's constant, spectral gap, network efficiency) are combined
    with stationary passenger throughput into a stop criticality score, with
    reducible chains handled explicitly.
  \item[C6.] \textbf{Robustness-driven augmentation.}  The node-removal
    analysis is inverted into a link-augmentation problem.  An exact
    shortest-path update makes greedy search over all candidate links cheap,
    three structural propositions explain which links can and cannot help, and
    the suggested links are compared with the projects planned for Budapest.
\end{description}

\subsection{Paper Organisation}

Section~\ref{sec:network} formalises the network graph.
Section~\ref{sec:entropy} derives the entropy-maximising flow and transition
probabilities.
Section~\ref{sec:ctmc} develops the CTMC generator and the Kolmogorov Forward
Equation.
Section~\ref{sec:prior} presents the hierarchical demand prior.
Section~\ref{sec:ssa} covers the Gillespie SSA and $\tau$-leaping.
Section~\ref{sec:complexity} gives a computational complexity analysis.
Section~\ref{sec:impl} describes the Python implementation.
Section~\ref{sec:results} reports empirical results on the full BKK network.
Section~\ref{sec:resilience} develops the resilience analysis.
Section~\ref{sec:augment} uses it to suggest new links for a more robust
network.
Section~\ref{sec:limitations} discusses limitations and future work.
Section~\ref{sec:conclusion} concludes.

% =============================================================================
\section{Network Graph and GTFS Data Model}
\label{sec:network}

\subsection{GTFS Inventory}

\begin{table}[htbp]
\centering
\small
\begin{tabular}{lll}
\toprule
GTFS File & Key Fields Used & Role in Model \\
\midrule
\texttt{stops.txt}       & \texttt{stop\_id, stop\_lat, stop\_lon}     & Vertex set $S$, $|S|=n$ \\
\texttt{routes.txt}      & \texttt{route\_id, route\_type}             & Mode classification $v$ \\
\texttt{trips.txt}       & \texttt{trip\_id, route\_id, service\_id}   & Trip--route link \\
\texttt{stop\_times.txt} & \texttt{arrival\_time, departure\_time}     & Scheduled travel times $\bar\tau^v_{ij}$ \\
\texttt{shapes.txt}      & \texttt{shape\_pt\_lat, shape\_pt\_lon}     & Route distance $d^v_{ij}$ \\
\texttt{calendar.txt}    & \texttt{monday\ldots friday}                & Weekday service filter \\
\texttt{transfers.txt}   & \texttt{from\_stop\_id, min\_transfer\_time}& Transfer penalty $\delta_i$ \\
\bottomrule
\end{tabular}
\caption{BKK GTFS files used in this work.}
\label{tab:gtfs}
\forras{BKK GTFS Schedule feed (\url{https://opendata.bkk.hu}, CC0-1.0)}
\end{table}

\subsection{Multi-Modal Stop Graph}

\begin{definition}[Multi-modal stop graph]
Let $S=\{1,\ldots,n\}$ denote the set of GTFS stops and
$\mathcal{V}=\{1,\ldots,m\}$ the set of vehicle modes.  For each mode
$v\in\mathcal{V}$, define the modal subgraph $G^v=(S^v,E^v)$, where
$S^v\subseteq S$ is the set of stops served by mode $v$ and
\[
  E^v=\bigl\{(i,j)\in S^v\times S^v:\exists\text{ a trip of mode }v
    \text{ visiting }i\text{ immediately before }j\bigr\}.
\]
The multi-modal graph is the union $G=(S,\bigcup_v E^v)$.
\end{definition}

\begin{remark}[Sparsity]
On the BKK network $|E^v|\ll|S^v|^2$ for all modes.  The bus sub-graph has
a fill ratio $|E^v|/|S^v|^2$ of about $0.03\%$ (Table~\ref{tab:network}).  All matrices are stored in CSR format; the full
$n\times n$ dense matrix is never materialised.
\end{remark}

\begin{definition}[Neighbourhood sets]
$\mathcal{N}^v(i)=\{j\in S^v:(i,j)\in E^v\}$,\quad
$\mathcal{N}^v_{\mathrm{in}}(i)=\{j\in S^v:(j,i)\in E^v\}$.
\end{definition}

\subsection{BKK Mode Classification}

\begin{table}[htbp]
\centering
\small
\begin{tabular}{cllcccl}
\toprule
$v$ & Mode & \texttt{rt} & Speed (km/h) & Cap.\ (pax) & $\kappa^v$ & Example Lines \\
\midrule
1 & Bus          & 3   & 18--25 & 80--180   & 1.00 & 7, 8E, 100E, 200E \\
2 & Tram         & 0   & 15--22 & 200--350  & 0.70 & 2, 4, 6, 47--49 \\
3 & Trolleybus   & 800 & 15--20 & 100--170  & 0.85 & 70, 72, 74, 76 \\
4 & Metro        & 1   & 30--40 & 800--1000 & 0.40 & M1, M2, M3, M4 \\
5 & Suburban rly & 2   & 25--40 & 400--600  & 0.55 & H5, H6, H7, H8 \\
6 & Ferry        & 4   & 12--18 & 100       & 1.20 & D11, D12, D14 \\
\bottomrule
\end{tabular}
\caption{BKK vehicle modes.  $\kappa^v$ is the mode-specific comfort weight;
  \texttt{rt} is the \texttt{route\_type} code recognised by the framework.}
\label{tab:modes}
\forras{BKK extended GTFS \texttt{route\_type} codes}
\end{table}

\subsection{Admissible Transitions and Disconnected Pairs}

A transition $(i,j)$ on mode $v$ is admissible iff $(i,j)\in E^v$.  We assign
a finite cost only to admissible transitions; all others receive
$C^v_{ij}=+\infty$, yielding $\exp(-\lambda^v C^v_{ij})=0$ and zero
probability/rate.  This handles graph sparsity exactly: we do
not approximate disconnected pairs with a large finite penalty~$C_\infty$,
which would inject spuriously small but non-zero probabilities across the entire
network.

\begin{remark}[Self-loops]
The entropy-maximisation formulation excludes self-loops ($i\ne j$ in the
constraint set); accordingly $C^v_{ii}$ plays no role.  The correct
correspondence is that $P^v$ as defined in Definition~\ref{def:dtmc} is the
\emph{embedded jump chain} of the CTMC, not the $\Delta t$-uniformisation.
\end{remark}

\subsection{Station-Level Aggregation of Metro Platforms}
\label{sec:platforms}

The BKK feed represents each metro station by two or more platform-level
\texttt{stop\_id} values, one per travel direction, linked to a common
\texttt{parent\_station}.  Built naively, the metro graph therefore consists of
one-directional paths: every trip enters a platform and leaves from the same
platform in the same direction, so no platform can be reached again once left.
Such a graph is a directed acyclic graph, it has no non-trivial strongly
connected component, and the corresponding chain has no stationary
distribution with support on the metro network.

We therefore define a node map $\sigma:S\to\tilde S$,
\begin{equation}
  \sigma(s)=\begin{cases}
    \texttt{parent\_station}(s), & \text{if defined},\\
    s, & \text{otherwise},
  \end{cases}
\end{equation}
and build the modal graphs on $\tilde S$.  Edges with $\sigma(i)=\sigma(j)$ are
dropped, travel times and distances of parallel edges are aggregated by the
median, and the departure hazard of a station is computed from the number of
distinct peak trips departing from any of its platforms.  In the feed used here
537 \texttt{stops.txt} records have a parent station (platforms, entrances and
generic nodes).  Among the stops actually served, the 103 metro platforms
collapse into 48 stations, and the resulting metro graph (48 nodes, 96 directed
edges) is strongly connected, so its chain is irreducible.  Only 26 served
surface platforms (22 bus, 4 tram) have a parent station; all other surface
stops keep their platform-level identifiers.

% =============================================================================
\section{Entropy Maximisation and the Softmax Flow}
\label{sec:entropy}

\subsection{Generalised Cost Matrix}

\begin{definition}[Generalised cost]
For admissible pair $(i,j)\in E^v$,
\begin{equation}
  C^v_{ij} = \kappa^v\,\bar\tau^v_{ij} + \gamma\,d^v_{ij}
              + \delta\,\mathbf{1}_{\{i\in\mathcal{T}\}}\,t_{\mathrm{tf}}(i),
  \label{eq:cost}
\end{equation}
where $\bar\tau^v_{ij}$~[s] is the median scheduled travel time,
$\kappa^v>0$ is the mode-specific comfort weight,
$d^v_{ij}$~[m] is the in-vehicle distance,
$\gamma$~[s\,m$^{-1}$] is the distance penalty,
$t_{\mathrm{tf}}(i)$~[s] is the minimum transfer time at stop~$i$,
and $\delta\in\{0,1\}$ indicates a mode change.  All terms are in seconds.
\end{definition}

\begin{remark}[Dimensional consistency]
$C^v_{ij}$~[s] and $\lambda^v$~[s$^{-1}$] make the exponent
$\lambda^v C^v_{ij}$ dimensionless, consistent with the interpretation of
$\lambda^v$ as an inverse temperature governing cost sensitivity.
\end{remark}

\subsection{Entropy-Maximising Flow}

\begin{assumption}[Separable origin constraint]
\label{ass:origin}
Passengers are assigned a departure mode~$v$ before departure; within a mode,
they choose destinations to maximise entropy subject to an expected-cost
constraint.
\end{assumption}

\begin{proposition}[Wilson softmax flow]
Under Assumption~\ref{ass:origin}, the flow $F^v_{ij}\ge0$ maximising Shannon
entropy $H=-\sum_{v,i,j}F^v_{ij}\ln F^v_{ij}$ subject to the origin
constraint $\sum_j F^v_{ij}=N^v_i(0)$ and the cost constraint
$\sum_{ij}F^v_{ij}C^v_{ij}\le\bar{C}^v$ is given by
\begin{equation}
  F^v_{ij} = N^v_i(0)\,\pi^v_{j|i},\qquad
  \pi^v_{j|i} = \frac{\exp(-\lambda^v C^v_{ij})}
    {\displaystyle\sum_{k\in\mathcal{N}^v(i)}\exp(-\lambda^v C^v_{ik})},
  \label{eq:softmax}
\end{equation}
where $\lambda^v\ge0$ is the Lagrange multiplier, uniquely determined by the
binding cost constraint.
\end{proposition}

\begin{proof}[Proof sketch]
The Lagrangian yields $\ln F^v_{ij}=-1+\mu^v_i-\lambda^v C^v_{ij}$, i.e.\
$F^v_{ij}\propto\exp(-\lambda^v C^v_{ij})$.  Row normalisation gives
\eqref{eq:softmax}.  Convexity of $-H$ ensures a unique maximiser.
\end{proof}

\begin{definition}[DTMC transition matrix]
\label{def:dtmc}
The row-stochastic DTMC transition matrix for mode~$v$ is
$P^v_{ij}=\pi^v_{j|i}$ ($i\ne j$), $P^v_{ii}=0$ (no self-loop),
$P^v_{ij}=0$ for $(i,j)\notin E^v$.
\end{definition}

\subsection{DTMC Time Evolution}

The population vector on mode~$v$ at step~$t$ evolves as
\begin{equation}
  N^v(t+1) = N^v(t)\,P^v,\qquad N_i(t)=\sum_{v=1}^m N^v_i(t).
\end{equation}

% =============================================================================
\section{Continuous-Time Markov Chain and the Kolmogorov Forward Equation}
\label{sec:ctmc}

\subsection{Motivation for the CTMC Formulation}

The GTFS records precise departure timestamps at each stop.  Rather than
discretising the schedule into fixed steps $\Delta t$, we embed the passenger
movement process in continuous time.  Under the memoryless (Markov)
approximation for inter-departure intervals---justified when vehicles run at
roughly constant headway---the waiting time is exponentially distributed and the
process is a CTMC.

\begin{assumption}[Poisson arrivals]
\label{ass:poisson}
The departure process at stop~$i$ on mode~$v$ is approximated as a Poisson
process with rate $f^v_i$ [vehicles\,s$^{-1}$] estimated from the GTFS
timetable.
\end{assumption}

\subsection{CTMC Generator}

\begin{definition}[Infinitesimal generator]
For mode~$v$, the generator $Q^v\in\mathbb{R}^{n\times n}$ is defined by
\begin{equation}
  Q^v_{ij} = \begin{cases}
    q^v_{ij}\ge0, & j\ne i,\ (i,j)\in E^v, \\
    0,             & j\ne i,\ (i,j)\notin E^v, \\
    -\displaystyle\sum_{k\in\mathcal{N}^v(i)} q^v_{ik}, & j=i.
  \end{cases}
\end{equation}
All off-diagonal entries are non-negative and all row sums vanish.
\end{definition}

\subsection{Rates from Entropy Kernel and GTFS Frequencies}
\label{sec:rates}

The off-diagonal rate decomposes as
\begin{equation}
  q^v_{ij} =
  \underbrace{\mu^v_i}_{\text{departure intensity}}\cdot
  \underbrace{\pi^v_{j|i}}_{\text{Wilson kernel}},
  \qquad (i,j)\in E^v,\;j\ne i,
  \label{eq:rate}
\end{equation}
where the per-passenger departure hazard is the scheduled vehicle frequency
\begin{equation}
  \mu^v_i = f^v_i
  \quad[\text{s}^{-1}].
  \label{eq:intensity}
\end{equation}
Here $f^v_i$ [vehicles\,s$^{-1}$] is the peak-window departure frequency.
This choice is dictated by dimensional consistency: every entry of a CTMC
generator has units s$^{-1}$.  Capacity $\bar{k}^v$ and load factor $\rho^v$ describe
supply and demand, but multiplying them into $Q$ and then again by $N_i$ in
the KFE/SSA propensity would produce the invalid unit pax$^2$\,s$^{-1}$.

\subsection{DTMC--CTMC Equivalence}

\begin{proposition}[Matrix exponential embedding]
The matrix exponential $P^v(\Delta t):=\exp(Q^v\Delta t)$ is stochastic for
all $\Delta t\ge0$; its first-order expansion recovers the Wilson DTMC entry
structure with rescaled self-loop probabilities.  The DTMC of
Definition~\ref{def:dtmc} is the \emph{embedded jump chain} of the CTMC,
observed at transition moments stripped of holding times.
\end{proposition}

\subsection{Kolmogorov Forward Equation}

\begin{theorem}[KFE]
Let $p^v(t)\in[0,1]^n$ denote the normalised probability distribution of a
single tagged passenger of mode~$v$.  Then
\begin{equation}
  \frac{\mathrm{d}p^v(t)}{\mathrm{d}t} = p^v(t)\,Q^v,\qquad
  p^v(0)=p^v_0,\label{eq:kfe}
\end{equation}
with unique solution $p^v(t)=p^v_0\exp(Q^v t)$.
\end{theorem}

By linearity and the mean-field (fluid-limit) approximation, the same equation
holds for population counts:
\begin{equation}
  \frac{\mathrm{d}N^v(t)}{\mathrm{d}t} = N^v(t)\,Q^v,\qquad
  N^v(0)=N^v_0.\label{eq:kfe-pop}
\end{equation}
Conservation of mass follows from $Q^v\mathbf{1}=\mathbf{0}$.

\subsection{Stationarity, Ergodicity, and Existence of Stationary Distributions}

\begin{assumption}[Irreducibility per mode]
\label{ass:irreducible}
For each mode~$v$, the directed graph~$G^v$ is strongly connected.
\end{assumption}

\begin{proposition}[Ergodicity]
Under Assumption~\ref{ass:irreducible}, $Q^v$ is irreducible and the CTMC is
ergodic.  A unique stationary distribution $\pi^v$ with $\pi^v Q^v=\mathbf{0}$,
$\|\pi^v\|_1=1$ exists, and $p^v(t)\to\pi^v$ exponentially fast.
\end{proposition}

\begin{remark}[Irreducibility in practice]
The BKK graph is not strongly connected for every mode in isolation: the
tram and bus graphs have 97 and 492 strongly connected components
(\texttt{scipy.sparse.csgraph}).  For such chains the stationary distribution
is not unique; we use the uniform-start Ces\`aro average and treat
irreducibility-dependent quantities explicitly (Section~\ref{sec:wellposed}).
\end{remark}

\begin{remark}[Time-homogeneity and GTFS schedules]
A real GTFS schedule is time-varying; the time-homogeneous CTMC is valid only
within a narrow service window (e.g., 07:00--09:00 peak hour) where $f^v_i$
is approximately constant.
\end{remark}

% =============================================================================
\section{Prior Estimation of Initial Passenger Counts}
\label{sec:prior}

\subsection{Setting and Identifiability}

In the absence of APC or AFC data, the initial count $N_i(0)$ is unobserved.
We construct a prior $\hat{N}_i(0)$ from freely available covariates and
calibrate its overall scale to BKK's published aggregate ridership.

\begin{assumption}[Log-linear demand]
\begin{equation}
  N_i(0)\sim\mathrm{Poisson}(\lambda_i),\qquad
  \log\lambda_i = \beta_0+\boldsymbol{\beta}^\top\mathbf{x}_i+\xi_i,\quad
  \xi_i\overset{\mathrm{iid}}{\sim}\mathcal{N}(0,\sigma^2).
  \label{eq:loglinear}
\end{equation}
\end{assumption}

\subsection{Feature Vector}

\begin{table}[htbp]
\centering
\small
\begin{tabular}{lllll}
\toprule
Feature & Symbol & Source & Unit \\
\midrule
Peak-hour service frequency     & $f_i$       & GTFS             & trips\,h$^{-1}$ \\
Population in 400--800\,m catchment & $\mathrm{Pop}_i$ & KSH/WorldPop & persons \\
POI count in 400\,m catchment   & $\mathrm{POI}_i$ & OpenStreetMap  & count \\
Network betweenness centrality  & $B_i$       & GTFS graph       & [0,1] \\
Interchange indicator           & $I_i$       & GTFS transfers   & $\{0,1\}$ \\
Distance to De\'ak F.\ t\'er   & $D_i$       & Haversine        & km \\
\bottomrule
\end{tabular}
\caption{Features used in the log-linear demand prior.}
\label{tab:features}
\forras{own compilation from the transit-demand literature}
\end{table}

\subsection{Literature-Informed Prior Coefficients}

\begin{table}[htbp]
\centering
\small
\begin{tabular}{lcccc}
\toprule
Covariate & Symbol & $\mu_\beta$ & $\sigma_\beta$ & Basis \\
\midrule
$\log(1+f_i)$                & $\beta_1$ & $+0.70$ & 0.10 & Service elasticity \citep{schiewe2022} \\
$\log(1+\mathrm{Pop}_i)$    & $\beta_2$ & $+0.60$ & 0.15 & Population elasticity \\
$\log(1+\mathrm{POI}_i)$    & $\beta_3$ & $+0.20$ & 0.10 & Activity generation \\
$B_i$                        & $\beta_4$ & $+0.30$ & 0.15 & Network centrality \\
$I_i$                        & $\beta_5$ & $+0.50$ & 0.20 & Transfer-hub effect \\
$\log(1+D_i)$                & $\beta_6$ & $-0.25$ & 0.10 & Distance decay \\
\bottomrule
\end{tabular}
\caption{Prior distributions $\mathcal{N}(\mu_\beta,\sigma_\beta^2)$ for
  log-linear coefficients.}
\label{tab:priors}
\forras{own compilation from the transit-demand literature}
\end{table}

\subsection{Calibration Anchor}

BKK reports $N_{\mathrm{day}}\approx4\times10^6$ weekday boardings.
With $\phi_{\mathrm{peak}}\approx0.11$,
\begin{equation}
  N_{\mathrm{peak}}=\phi_{\mathrm{peak}}\,N_{\mathrm{day}}\approx440{,}000\text{ passengers.}
\end{equation}
The intercept $\beta_0$ is set so that $\sum_i\hat{N}_i(0)=N_{\mathrm{peak}}$.

\subsection{Diurnal Profile}

$\hat{N}_i(0;h)=24\,\phi(h)\,\hat{N}_i(0)$ with
$\phi(07\text{:}00\text{--}09\text{:}00)\approx0.11$.

\subsection{Three Estimators}

\textbf{E1~--- Service proxy (GTFS only).}
$\hat{N}_i(0)\propto f_i\cdot\bar{k}^{v_i}\cdot\rho$.
No external data needed; captures frequency variation but ignores land use.

\textbf{E2~--- Poisson log-linear (recommended).}
Equation~\eqref{eq:loglinear} with priors from Table~\ref{tab:priors}.

\textbf{E3~--- Full Bayesian (for future use with AFC/APC data).}
Full posterior inference via MCMC.  Reduces to E2 when no observations are
available.

\subsection{Uncertainty Propagation}

Monte Carlo propagation: draw
$\boldsymbol{\beta}^{(s)}\sim\mathcal{N}(\boldsymbol{\mu}_\beta,\Sigma_\beta)$,
compute $\hat{N}^{(s)}_i(0)$, run KFE or SSA, report empirical quantiles.

% =============================================================================
\section{Gillespie Stochastic Simulation Algorithm}
\label{sec:ssa}

\subsection{Reaction-Channel Representation}

The propensity of channel $(v,i\to j)$ is
\begin{equation}
  a^v_{ij}(t)=N^v_i(t)\,q^v_{ij},\qquad
  a_0(t)=\sum_v\sum_i\sum_{j\in\mathcal{N}^v(i)}a^v_{ij}(t).
\end{equation}
This is exact for $N^v_i(t)\in\mathbb{Z}_{\ge0}$; the mean-field ODE
\eqref{eq:kfe-pop} is recovered in the limit $N^v_i\to\infty$.

\subsection{Direct Method (Exact SSA)}

Channel selection uses the $O(R)$ inverse-CDF method (\texttt{np.cumsum} and
\texttt{np.searchsorted}).  The Vose alias method \citep{vose1991} would give
$O(1)$ selection but must be rebuilt after every event, because each event
changes two propensities; for the channel counts considered here the simpler
inverse-CDF method is sufficient.

\begin{figure}[H]
\centering
\fbox{\begin{minipage}{0.92\textwidth}
\textbf{Algorithm 1: Gillespie Direct Method for BKK Multi-Modal CTMC}\\[4pt]
\textit{Require:} Sparse channels $\{(v,i,j,q^v_{ij})\}$ (only $(i,j)\in E^v$);
initial state $N(0)$; horizon $T$.\\
\textit{Ensure:} Event list $\mathcal{E}$; final state $N(T)$.\\[4pt]
\begin{enumerate}
  \item $t\leftarrow0$,\quad $\mathcal{E}\leftarrow\emptyset$
  \item \textbf{while} $t<T$ \textbf{do}
  \begin{enumerate}
    \item Compute propensities (vectorised):\quad
          $a[r]\leftarrow N^v_i\cdot q^v_{ij}$;\quad
          $a_0\leftarrow\sum_r a[r]$
    \item \textbf{if} $a_0=0$ \textbf{then break}
    \item Draw $\tau\leftarrow-\ln(u_1)/a_0$,\quad $u_1\sim\mathcal{U}(0,1)$
    \item \textbf{if} $t+\tau>T$ \textbf{then break}
    \item Select $r^*$ via inverse-CDF on $\{a[r]/a_0\}$ using $u_2\sim\mathcal{U}(0,1)$
    \item $t\leftarrow t+\tau$;\quad
          $N^{v^*}_{i^*}\leftarrow N^{v^*}_{i^*}-1$;\quad
          $N^{v^*}_{j^*}\leftarrow N^{v^*}_{j^*}+1$
    \item Append $(t,v^*,i^*,j^*)$ to $\mathcal{E}$
  \end{enumerate}
  \item \textbf{return} $\mathcal{E}$, $N(T)$
\end{enumerate}
\end{minipage}}
\caption{Gillespie Direct Method for BKK Multi-Modal CTMC.}
\label{alg:ssa}
\forras{own computation, \texttt{bkk/} framework}
\end{figure}

\subsection{$\tau$-Leaping Approximation}

When $N^v_i\gg1$, the $\tau$-leap \citep{gillespie2001} advances by a fixed
step~$\tau$ and draws
\begin{equation}
  \Delta k^v_{ij}\sim\mathrm{Poisson}(a^v_{ij}\tau),
  \label{eq:tauleap}
\end{equation}
then updates populations subject to an integer non-negativity guard.  The
implemented multinomial oversubscription rule preserves both mass and integer
counts, while remaining an approximate leap rather than exact SSA.

\textbf{Integer atomic outflow handling.} If channels sharing origin $i$
propose a total drain above $N_i^v$, independent serial clamping is
order-dependent and proportional scaling creates fractional passengers.  The
implementation therefore redraws the $N_i^v$ available passengers from a
multinomial distribution with probabilities proportional to proposed counts:
\begin{align}
  \Delta k^v_{ij} &\sim \mathrm{Poisson}(a^v_{ij}\tau), \label{eq:draw}\\
  D_{vi} &= \textstyle\sum_{j\in\mathcal{N}^v(i)}\Delta k^v_{ij},
    \label{eq:total-out}\\
  (\Delta k^v_{ij})_j \mid D_{vi}>N_i^v
  &\sim \operatorname{Multinomial}\!\left(
  N_i^v,\;(\Delta k^v_{ij}/D_{vi})_j\right).
\end{align}
Initial non-integer demand estimates are converted once by largest-remainder
rounding, preserving the rounded grand total.  Thereafter SSA and tau-leap
states remain integer-valued.

\textbf{Step-size control.}
The implementation uses the bounded propensity-ratio heuristic
\begin{equation}
  \tau = \operatorname{clip}\!\left(
    \varepsilon a_0/\max_r a_r,\;0.1\tau_0,\;10\tau_0\right).
\end{equation}
This is not the full Cao--Gillespie--Petzold method \citep{cao2006}, which
requires species-wise drift and variance bounds.

\subsection{Hybrid SSA--Fluid Decomposition}

At stops with $N^v_i>N_{\mathrm{thresh}}\approx1000$, the coefficient of
variation $\sim N_i^{-1/2}\ll1$ and the fluid ODE~\eqref{eq:kfe-pop} is
accurate; at stops with $N^v_i\lesssim10$ stochastic fluctuations dominate.
We recommend a hybrid regime with coupling via boundary propensities at
interface stops.

% =============================================================================
\section{Computational Complexity Analysis}
\label{sec:complexity}

\subsection{Graph-Structure Complexity}

Let $n=|S|$, $m=|\mathcal{V}|$, $R=\sum_v|E^v|$.  For the BKK feed used in
this paper (Section~\ref{sec:results}): $n=5{,}177$, $m=4$, $R=6{,}450$.

\subsection{One-Time Setup}

\begin{table}[htbp]
\centering
\small
\begin{tabular}{lll}
\toprule
Operation & Complexity & Notes \\
\midrule
Parse GTFS files                 & $O(|\mathrm{trips}|+|\mathrm{stop\_times}|)$ & pandas I/O \\
Build cost matrices $C^v$ (CSR) & $O(R)$                                       & Admissible pairs only \\
Build generator $Q^v$ (CSR)     & $O(R)$                                       & \\
Strongly-connected components   & $O(n+R)$                                     & Tarjan \\
Betweenness centrality          & $O(nm)$ sparse                               & NetworkX \\
\bottomrule
\end{tabular}
\caption{One-time setup complexity.}
\forras{own computation, \texttt{bkk/} framework}
\end{table}

\subsection{Simulation Complexity}

\textbf{KFE (ODE solver):} $O(R\,T_{\mathrm{steps}})$.

\textbf{Matrix exponential (Krylov):}
$O(R\ell)$ per time-point, $\ell=30$ default.

\textbf{Gillespie SSA:} $O(a_0 T)$ events, each costing $O(R)$ with
inverse-CDF selection, where $a_0=\sum_{v,i}N_i^v f_i^v$.  With the E1 prior
anchored to $N_{\mathrm{peak}}=440{,}000$ passengers the full network has
$a_0\approx5{,}600$\,s$^{-1}$, i.e.\ about $5.1\times10^6$ events in a 900\,s
window and ${\sim}3\times10^{10}$ elementary operations.  Exact SSA is
therefore used on the metro sub-network ($a_0\approx814$\,s$^{-1}$), and
$\tau$-leaping on the full network.

\textbf{$\tau$-leaping:} $O(R)$ Poisson draws per step (vectorised).
With a fixed step $\tau=1$\,s and $T=900$\,s this is at most
$900\times R\approx5.8\times10^6$ draws; the adaptive step-size rule needed
about 91 steps on the full network.

\textbf{Resilience:} Naive dense perturbation is $O(Kn^3)$.  ARPACK
truncation reduces the spectral work, while node deletion still materialises a
dense matrix; rank-one (Woodbury-type) updates \citep{hunter2014} are left for
future work.

\subsection{Memory Requirements}

\begin{table}[htbp]
\centering
\small
\begin{tabular}{lcc}
\toprule
Object & Dense & Sparse (CSR) \\
\midrule
Bus generator $Q^{\mathrm{bus}}$ & $\sim164$\,MB & $\sim0.14$\,MB \\
All $m$ generators                & $\sim167$\,MB & $\sim0.16$\,MB \\
Simulation state $(N^v_i)_{v,i}$  & $\sim166$\,kB & --- \\
\bottomrule
\end{tabular}
\caption{Memory footprint at BKK scale ($n=5{,}177$, $R=6{,}450$, $m=4$).}
\forras{own computation, \texttt{bkk/} framework}
\end{table}

The dense representation is three orders of magnitude larger than CSR; for
the node-removal sweep, which forms many perturbed matrices, sparse storage is
what keeps the analysis on a single workstation.

% =============================================================================
\section{Python Implementation}
\label{sec:impl}

\subsection{Software Stack}

\textbf{pandas / zipfile:} GTFS ingestion.\quad
\textbf{numpy / scipy.sparse:} sparse matrices (CSR throughout).\quad
\textbf{scipy.sparse.linalg.expm\_multiply:} Krylov action of $\exp(Qt)$.\quad
\textbf{scipy.integrate.solve\_ivp:} ODE integration of \eqref{eq:kfe-pop}.\quad
\textbf{networkx:} betweenness centrality, SCC decomposition.\quad
\textbf{overpy / osmnx:} OpenStreetMap POI queries.\quad
\textbf{rasterio:} WorldPop rasters for catchment population.

\subsection{Core Implementation}

\textbf{Listing 1: Generator construction and KFE solve.}
The \texttt{GeneratorBuilder.build()} method groups edges by origin stop using
\texttt{numpy} argsort, applies the log-sum-exp stabilised softmax within each
group (no Python-level for-loop over individual stop pairs), and assembles a
CSR generator.

\begin{lstlisting}[caption={Generator construction and KFE solve
  (\texttt{bkk/generator.py}, abridged).},label={lst:generator}]
import numpy as np
from scipy.sparse import lil_matrix, csr_matrix
from scipy.integrate import solve_ivp

def build(mg, lam: float) -> csr_matrix:
    """Build CTMC generator Q^v (CSR) without dense materialisation."""
    i_arr, j_arr, cost_arr, mu_arr = (
        mg.i_arr, mg.j_arr, mg.cost_arr, mg.mu_arr
    )
    n = mg.n_stops

    # Sort edges by origin stop for numpy groupby
    order = np.argsort(i_arr, kind='stable')
    i_s, j_s, c_s = i_arr[order], j_arr[order], cost_arr[order]
    q_vals = np.empty(len(order), dtype=np.float64)

    boundaries = np.where(np.diff(i_s))[0] + 1
    groups = np.split(np.arange(len(order)), boundaries)

    for grp_idx in groups:
        orig  = i_s[grp_idx[0]]
        costs = c_s[grp_idx]
        # Log-sum-exp stabilised softmax: pi^v_{j|i}
        log_w  = -lam * costs
        log_w -= log_w.max()        # numerical stability shift
        w = np.exp(log_w)
        pi = w / w.sum()
        # q^v_{ij} = mu^v_i * pi^v_{j|i}
        q_vals[grp_idx] = mu_arr[orig] * pi

    # Assemble CSR generator Q
    Q_lil = lil_matrix((n, n), dtype=np.float64)
    for r in range(len(i_s)):
        Q_lil[i_s[r], j_s[r]] = q_vals[r]
    for i in range(n):
        Q_lil[i, i] = -Q_lil.getrow(i).sum()  # row sum = 0
    return Q_lil.tocsr()


def kfe_solve_ode(Q: csr_matrix, N0: np.ndarray,
                  T: float, t_eval: np.ndarray) -> np.ndarray:
    """Solve mean-field ODE dN/dt = N Q via LSODA (sparse)."""
    def rhs(t, N):
        return Q.T.dot(N)      # column form: d(N^T)/dt = Q^T N^T
    sol = solve_ivp(rhs, [0, T], N0.astype(float),
                    t_eval=t_eval, method='LSODA',
                    rtol=1e-8, atol=1e-10)
    return sol.y.T             # shape (len(t_eval), n)
\end{lstlisting}

\textbf{Listing 2: Gillespie SSA --- Direct Method with cumsum channel
selection.}
The implementation matches Algorithm~\ref{alg:ssa} exactly, using
\texttt{np.searchsorted} on the cumulative propensity vector for $O(R)$ channel
selection.

\begin{lstlisting}[caption={Gillespie SSA (\texttt{bkk/simulate.py},
  inner loop).},label={lst:ssa}]
import numpy as np

def run_ssa(v_arr, i_arr, j_arr, q_arr, N, T, max_events, rng):
    """
    Exact Doob-Gillespie Direct Method (O(R) per event).

    v_arr, i_arr, j_arr : int arrays shape (R,) -- channel indices
    q_arr               : float array shape (R,) -- base rates q^v_{ij}
    N                   : float array shape (m, n) -- populations
    T, max_events, rng  : horizon, safety cap, RNG
    """
    t = 0.0
    events = []
    for _ in range(max_events):
        # Propensities: a[r] = N[v, i] * q[r]
        a  = N[v_arr, i_arr] * q_arr     # shape (R,)
        a0 = a.sum()
        if a0 <= 0.0:
            break

        # Waiting time ~ Exp(a0)
        tau = -np.log(rng.random()) / a0
        if t + tau >= T:
            break
        t += tau

        # Channel selection via inverse-CDF (O(R))
        u = rng.random() * a0
        r = int(np.searchsorted(np.cumsum(a), u))
        r = min(r, len(v_arr) - 1)

        v, ii, jj = v_arr[r], i_arr[r], j_arr[r]
        N[v, ii] = max(0.0, N[v, ii] - 1)
        N[v, jj] += 1
        events.append((t, v, ii, jj))

    return events, N
\end{lstlisting}

\textbf{Listing 3: integer $\tau$-leaping with atomic outflow handling.}
Oversubscribed origins use a multinomial allocation rather than fractional
proportional scaling.

\begin{lstlisting}[caption={$\tau$-leaping with atomic outflow clamping
  (\texttt{bkk/simulate.py}).},label={lst:tauleap}]
import numpy as np

def run_tauleap(v_arr, i_arr, j_arr, q_arr, N, T, tau, rng):
    """
    Integer tau-leaping with grouped multinomial oversubscription.
    """
    t = 0.0
    while t < T:
        tau_step = min(tau, T - t)

        # Propensities and Poisson draws for all R channels
        a  = N[v_arr, i_arr] * q_arr              # (R,)
        dk = rng.poisson(a * tau_step).astype(np.int64)

        # Sum proposed outflow per (v, i)
        total_out = np.zeros_like(N)              # (m, n)
        np.add.at(total_out, (v_arr, i_arr), dk)

        for v, i in np.argwhere(total_out > N):
            channels = np.flatnonzero((v_arr == v) & (i_arr == i))
            proposed = dk[channels]
            dk[channels] = rng.multinomial(
                int(N[v, i]), proposed / proposed.sum()
            )

        # Atomic update
        np.add.at(N, (v_arr, i_arr), -dk)
        np.add.at(N, (v_arr, j_arr), +dk)

        t += tau_step
    return N
\end{lstlisting}

\textbf{Listing 4: E2 prior estimator.}

\begin{lstlisting}[caption={E2 prior estimator for $\hat{N}_i(0)$
  (\texttt{bkk/demand.py}).},label={lst:prior}]
import numpy as np

BETA_PRIOR = np.array([0.70, 0.60, 0.20, 0.30, 0.50, -0.25])

def compute_features(freq_i, pop_i, poi_i, B_i, I_i, D_i):
    """Assemble log-transformed feature matrix X, shape (n, 6)."""
    return np.column_stack([
        np.log1p(freq_i),   # log(1 + f_i)
        np.log1p(pop_i),    # log(1 + Pop_i)
        np.log1p(poi_i),    # log(1 + POI_i)
        B_i,                # betweenness centrality [0, 1]
        I_i,                # interchange indicator {0, 1}
        np.log1p(D_i),      # log(1 + distance to Deak ter)
    ])

def estimate_N0(X: np.ndarray,
                beta: np.ndarray = BETA_PRIOR,
                N_peak: float = 440_000.0) -> np.ndarray:
    """Anchored prior: sum_i N_hat_i(0) = N_peak."""
    raw = np.exp(X @ beta)      # unnormalised intensities
    return raw * (N_peak / raw.sum())
\end{lstlisting}

% =============================================================================
\section{Empirical Results on the BKK Network}
\label{sec:results}

\subsection{Data and Network Summary}

The analysis uses the BKK GTFS Schedule feed released in June 2026
(\url{https://opendata.bkk.hu}, CC0-1.0).  The archive contains 6\,257 stop
records, 386 routes, 247\,616 trips and 678\,054 shape points.  The feed
encodes service days through \texttt{calendar\_dates.txt} only; selecting the
service of a Monday retains 99\,438 trips and 2\,040\,461 stop-time records.
Edges are built from all weekday trips; departure hazards $\mu^v_i$ from the
07:00--09:00 peak window (245\,866 departures).  After station-level
aggregation (Section~\ref{sec:platforms}) the network has $n=5{,}177$ nodes and
$R=6{,}450$ directed edges in four modes (Table~\ref{tab:network}).

\begin{table}[htbp]
\centering
\small
\begin{tabular}{llrrrrrr}
\toprule
Mode & \texttt{rt} & $|S^v|$ & $|E^v|$ & $|E^v|/|S^v|$ & Fill (\%) & $\mathrm{med}\,1/\mu^v_i$ (s) & $\mathrm{med}\,C^v_{ij}$ (s) \\
\midrule
Tram & 0 & 651 & 689 & 1.06 & 0.163 & 112 & 86 \\
Metro & 1 & 48 & 96 & 2.00 & 4.167 & 30 & 46 \\
Bus & 3 & 4\,525 & 5\,663 & 1.25 & 0.028 & 200 & 65 \\
Ferry & 4 & 2 & 2 & 1.00 & 50.000 & 900 & 218 \\
\midrule
\textbf{Total} & --- & 5\,177 & 6\,450 & --- & --- & --- & --- \\
\bottomrule
\end{tabular}
\caption{Network summary by mode (weekday service, 07:00--09:00 peak window).}
\label{tab:network}
\par\smallskip{\footnotesize Fill $=|E^v|/|S^v|^2$. $1/\mu^v_i$ is the scheduled headway at stops with at least one peak departure. The total stop count is over distinct nodes; a node served by several modes is counted once.}
\forras{own computation from the BKK GTFS Schedule feed (CC0-1.0)}
\end{table}


The graphs are extremely sparse: a surface stop has on average 1.06 (tram) or
1.25 (bus) scheduled successors, so for most stops the Wilson kernel has a
single admissible destination and route choice happens only at branching
stops.  The metro graph is different in kind: a station has on average two
successors (one per direction, more at interchanges, one at termini), the fill
ratio is two orders of magnitude higher, and the median headway is 30\,s,
against 112.5\,s for tram and 200\,s for bus stops.

\textbf{Mode coverage.}  The framework recognises trolleybus and suburban rail
under the extended \texttt{route\_type} codes 800 and 2 (Table~\ref{tab:modes}).
The current BKK feed codes them as 11 (16 trolleybus routes) and 109 (5 H\'EV
routes), so these two modes are not part of the network analysed below.

\subsection{Prior Estimate of Initial Counts}

The results use the E1 service-proxy estimator, $\hat N_i(0)\propto\mu_i$,
anchored to $N_{\mathrm{peak}}=440{,}000$; the E2 covariates (catchment
population, POI counts) require external rasters and are left to the
calibration stage.  The prior is moderately concentrated: its Gini coefficient
is 0.480, the 50 largest stops hold 5.6\% of all passengers, and the maximum,
1\,192 passengers, is at De\'ak Ferenc t\'er, where three metro lines meet
(Table~\ref{tab:prior-top10}).  The top of the list is formed by metro
interchanges and by the bus platforms that feed them (Keleti p\'alyaudvar,
Blaha Lujza t\'er, Astoria).

\begin{table}[htbp]
\centering
\small
\begin{tabular}{rlllrr}
\toprule
Rk & Stop & \texttt{stop\_id} & Mode & $f_i$ (dep.\,h$^{-1}$) & $\hat N_i(0)$ \\
\midrule
1 & De\'ak Ferenc t\'er & \texttt{CSF00954} & Metro & 333 & 1192 \\
2 & K\'alvin t\'er & \texttt{CS056227} & Metro & 242 & 868 \\
3 & Keleti p\'alyaudvar & \texttt{CS056233} & Metro & 241 & 863 \\
4 & Keleti p\'alyaudvar M & \texttt{F01159} & Bus & 200 & 714 \\
5 & Blaha Lujza t\'er M & \texttt{F01166} & Bus & 161 & 576 \\
6 & Keleti p\'alyaudvar M & \texttt{F01314} & Bus & 160 & 571 \\
7 & Zugl\'o vas\'ut\'allom\'as & \texttt{F02730} & Bus & 146 & 524 \\
8 & Zugl\'o vas\'ut\'allom\'as & \texttt{F02725} & Bus & 140 & 503 \\
9 & Ur\'ania & \texttt{F01171} & Bus & 136 & 487 \\
10 & Astoria M & \texttt{F01017} & Bus & 136 & 487 \\
\bottomrule
\end{tabular}
\caption{Top-10 stops by the E1 prior $\hat N_i(0)$, anchored to $N_{\mathrm{peak}}=440{,}000$.}
\label{tab:prior-top10}
\par\smallskip{\footnotesize Metro platforms are aggregated to their parent station; surface stops keep one \texttt{stop\_id} per platform, so a name can appear more than once.}
\forras{own computation from the BKK GTFS Schedule feed (CC0-1.0)}
\end{table}


\subsection{KFE Dynamics}

The population KFE~\eqref{eq:kfe-pop} was solved over $T=900$\,s with
$\lambda^v=1/300$\,s$^{-1}$ for all modes, using the Krylov action of
$\exp(Q^{v\top}t)$.  The largest relative change in total population over the
horizon is $1.4\times10^{-15}$, i.e.\ passengers are conserved to machine
precision, and all generators satisfy $\max_i|\sum_j Q^v_{ij}|<10^{-17}$ and
$\max_i|\sum_j P^v_{ij}-1|<3\times10^{-16}$.

\begin{table}[htbp]
\centering
\small
\begin{tabular}{lrrrrrr}
\toprule
Stop & $t=0$ & 60\,s & 120\,s & 300\,s & 600\,s & 900\,s \\
\midrule
De\'ak Ferenc t\'er & 1192 & 544 & 527 & 511 & 504 & 501 \\
K\'alvin t\'er & 868 & 519 & 504 & 492 & 485 & 480 \\
Keleti p\'alyaudvar & 863 & 415 & 405 & 388 & 370 & 360 \\
Keleti p\'alyaudvar M & 714 & 230 & 237 & 132 & 63 & 40 \\
Blaha Lujza t\'er M & 576 & 315 & 223 & 124 & 56 & 27 \\
\midrule
Sum, initial top-50 & 24\,517 & 19\,749 & 18\,054 & 16\,464 & 15\,590 & 15\,142 \\
Gini of $N(t)$ & 0.480 & 0.443 & 0.440 & 0.461 & 0.506 & 0.546 \\
\bottomrule
\end{tabular}
\caption{Mean-field (KFE) passenger counts at the five stops with the largest initial count, full network, $\lambda^v=1/300$\,s$^{-1}$.}
\label{tab:kfe-traj}
\forras{own computation from the BKK GTFS Schedule feed (CC0-1.0)}
\end{table}


Table~\ref{tab:kfe-traj} shows two time scales.  At metro stations the
population halves within the first minute (De\'ak Ferenc t\'er: 1\,192 to 544)
and then stays nearly constant: with a median headway of 30\,s the metro chain
mixes within a few headways, and what remains is its quasi-stationary level.
The high-frequency bus platforms (Keleti p\'alyaudvar M, Blaha Lujza t\'er M)
lose passengers just as quickly but do not level off; they fall to 40 and 27
passengers by $t=900$\,s.  The difference is structural: after aggregation the
metro station graph is traversable in both directions, so passengers
recirculate, whereas the bus graph is a collection of one-directional paths
along which passengers are carried downstream towards the termini.

The Gini coefficient of $N(t)$ first falls, from 0.480 to 0.440 at
$t=120$\,s, as the initial concentration at high-frequency stops spreads
along the lines, and then rises to 0.546 at $t=900$\,s.  The rise is driven
by \emph{absorbing} nodes: 469 of the 5\,226 mode-specific nodes either have no
peak-hour departure ($\mu^v_i=0$) or no outgoing edge, typically the arrival
platforms of terminal stops.  They hold 1.5\% of the passengers at $t=0$ and
19.2\% at $t=900$\,s.  The five largest stops at the horizon are all of this
kind (for example \'Ujpalota, Ny\'irpalota \'ut: 360 to 3\,473 passengers;
Sz\'ell K\'alm\'an t\'er M: 268 to 3\,001).  In an origin-constrained model
without trip ends this is how alighting at the end of a line appears, and it
is the main reason why the time-homogeneous KFE is meaningful only over a
window of the order of the typical trip duration.

\subsection{SSA and $\tau$-Leap Against the Mean Field}

For a linear population CTMC the KFE is the exact equation of the first moment,
so the comparison below is a test of the stochastic layers rather than of the
model.  At full scale an exact SSA run would require about $5\times10^6$
events per 900\,s (Section~\ref{sec:complexity}), so the exact comparison is
made on the metro sub-network, with the metro share of the prior (21\,251
passengers), $T=300$\,s and $M=20$ independent realisations.

\begin{table}[htbp]
\centering
\small
\begin{tabular}{lrrrrrr}
\toprule
Station & $\bar N^{\mathrm{SSA}}$ & $N^{05}$ & $N^{95}$ & KFE & CV (\%) & $N^{-1/2}$ (\%) \\
\midrule
Bajcsy-Zsilinszky \'ut & 574.0 & 545 & 609 & 567.0 & 4.3 & 4.2 \\
Pillang\'o utca & 545.0 & 512 & 584 & 550.5 & 4.6 & 4.3 \\
II. J\'anos P\'al p\'apa t\'er & 545.8 & 516 & 576 & 543.4 & 4.5 & 4.3 \\
Pusk\'as Ferenc Stadion & 517.4 & 497 & 539 & 520.5 & 3.0 & 4.4 \\
Blaha Lujza t\'er & 521.6 & 495 & 548 & 519.5 & 3.7 & 4.4 \\
R\'ak\'oczi t\'er & 522.4 & 489 & 556 & 516.7 & 4.1 & 4.4 \\
\bottomrule
\end{tabular}
\caption{Exact SSA ensemble ($M=20$) against the KFE mean on the metro sub-network at $T=300$\,s, six largest stations.}
\label{tab:ssa-kfe}
\par\smallskip{\footnotesize CV: ensemble coefficient of variation; $N^{-1/2}$: the Poisson reference $1/\sqrt{N^{\mathrm{KFE}}}$.}
\forras{own computation from the BKK GTFS Schedule feed (CC0-1.0)}
\end{table}


Every realisation completed the full horizon (on average 223\,571 events per
run), so no run is censored by the event budget.  Over all 48 stations the
ensemble mean deviates from the KFE solution by 0.84\% in relative $\ell^1$
norm, and 43 of the 48 stations (90\%) lie within two standard errors of the
KFE value, consistent with the nominal 95\% level for 48 stations.  The ensemble
coefficient of variation at the large stations is 3.0--4.6\%, close to the
reference value $N^{-1/2}$ (Table~\ref{tab:ssa-kfe}): at station level the
fluctuations are essentially those of a Poisson count, and the mean field is an
accurate description of every station with more than a few hundred passengers.

On the full network, ten $\tau$-leap realisations (initial step
$\tau_0=1$\,s, $\varepsilon=0.03$, on average 91 adaptive steps) conserve the
integer population exactly and reproduce the KFE state at $t=900$\,s with a
relative $\ell^1$ deviation of 2.4\%.  For each of the five largest stops at the
horizon the 5--95\% ensemble band contains the KFE value.  The deviation is
larger than for exact SSA, as expected of an approximate leap method with
multinomial redraws at nearly empty stops, but it remains small compared with
the uncertainty in the prior itself.

\subsection{Sensitivity to $\lambda$}

Sweeping $\lambda\in[10^{-5},10^{-1}]$\,s$^{-1}$ (25 log-spaced values) on
the bus sub-network gives the following picture.  (1)~The mean row entropy of
the transition matrix, $\bar H(P)=-\frac1n\sum_i\sum_jP_{ij}\ln P_{ij}$,
decreases monotonically from 0.166 to 0.077\,nats; its low level reflects the
fact that most bus stops have a single successor.  (2)~The Gini coefficient of
the (Ces\`aro) stationary distribution is essentially flat, between 0.53 and
0.54, for $\lambda\le5\times10^{-3}$\,s$^{-1}$, and then rises steeply to
0.81 at $\lambda=0.1$\,s$^{-1}$, when routing at branching stops becomes
nearly deterministic.  (3)~The default $\lambda=1/300$\,s$^{-1}$
($\bar H=0.163$, Gini 0.53) lies at the upper end of the flat regime, where the
results are insensitive to the exact value of $\lambda$; this matters because
$\lambda$ cannot be estimated without fare-collection data.

% =============================================================================
\section{Network Resilience Analysis}
\label{sec:resilience}

\subsection{Framework}

\textbf{Importance:}  the stationary departure flow of mode $v$ at stop $i$,
\begin{equation}
  I^v_i = N^v_{\mathrm{tot}}\,\pi^v_i\,\mu^v_i
  \quad [\text{pax\,s}^{-1}],
  \label{eq:importance}
\end{equation}
where $N^v_{\mathrm{tot}}$ is the prior population of mode $v$ and $\pi^v$ the
stationary (for reducible chains: uniform-start Ces\`aro) distribution.

\textbf{Vulnerability:}  Fractional degradation upon removing stop~$i^*$
(delete row/column $i^*$ from $P^v$, renormalise):
\begin{align}
  V^{\mathrm{Kem}}_{i^*} &=
  \frac{K(P^{(-i^*)})-K(P)}{|K(P)|},\\
  V^{\mathrm{gap}}_{i^*} &= \frac{g(P)-g(P^{(-i^*)})}{g(P)},\\
  V^{\mathrm{eff}}_{i^*} &= \frac{E(P)-E(P^{(-i^*)})}{E(P)},
\end{align}
where, only for an irreducible chain, Kemeny's constant is
$K(P)=\sum_{k=2}^n 1/(1-\lambda_k(P))$
\citep{kemeny1976,levene2002}, $g(P)=1-|\lambda_2(P)|$, and
$E(P)=\frac{1}{n(n-1)}\sum_{i\ne j}1/d_{ij}$ \citep{latora2001}, with $d_{ij}$
the hop distance in the support graph of $P$ (for $n>200$ estimated from the
first 200 source nodes).  $V^{\mathrm{Kem}}$, $V^{\mathrm{gap}}$ and
$V^{\mathrm{eff}}$ are truncated at zero.
Dead-end rows are made absorbing.  If node removal makes an irreducible chain
reducible, $V^{\mathrm{Kem}}=1$ is an explicit fragmentation penalty rather
than a fictitious Kemeny value.  For reducible baselines, $K$ is reported as
undefined and a uniform-start Ces\`aro stationary distribution is used.

\textbf{Criticality:}
\begin{equation}
  C_i = \sqrt{\tilde{I}_i\,\tilde{V}_i},\qquad
  \tilde{V}_i = \frac{1}{|\mathcal{M}^v|}\sum_{c\in\mathcal{M}^v}\tilde{V}^{c}_i,
  \label{eq:criticality}
\end{equation}
where every term is min-max normalised to $[0,1]$ over the candidate set of
mode $v$, and $\mathcal{M}^v\subseteq\{\mathrm{Kem},\mathrm{gap},\mathrm{eff}\}$
is the set of components that are defined for the baseline chain of mode $v$:
$V^{\mathrm{Kem}}$ requires an irreducible chain and $V^{\mathrm{gap}}$ a
positive spectral gap $g(P)>0$.  Because of the normalisation, $C_i$ ranks
stops within a mode; it is not comparable across modes.

\subsection{Well-Posedness of the Components on the BKK Chains}
\label{sec:wellposed}

On the BKK network the spectral-gap component is undefined for every mode, for
two different reasons.  The tram and bus chains are reducible (97 and 492
strongly connected components, the largest with 234 and 4\,009 nodes), so the
eigenvalue $1$ is multiple and $g(P)=0$.  The metro chain is irreducible but
periodic with period~2: every line is traversed in both directions, so the
station graph is bipartite and $-1$ is an eigenvalue of $P^{\mathrm{metro}}$
(the next largest moduli are $0.987$).  Again $g(P)=1-|\lambda_2|=0$, and the
relative change $(g-g^{(-i)})/g$ is a ratio of two numerically zero quantities.
We therefore use $\mathcal{M}^{\mathrm{metro}}=\{\mathrm{Kem},\mathrm{eff}\}$
and $\mathcal{M}^{\mathrm{tram}}=\mathcal{M}^{\mathrm{bus}}=\{\mathrm{eff}\}$.
Kemeny's constant is well defined for the periodic metro chain
($K=302.2$), since it only requires $\lambda_k\ne1$ for $k\ge2$.  The
two-node ferry chain is not ranked.

\subsection{Computational Acceleration}

Two speedups are used: (a)~node removal is evaluated for the top-$K=100$
candidates by importance in each mode (all 48 stations for the metro), and
(b)~for $n>50$ only the largest-modulus eigenvalues are computed with ARPACK;
for $n\le50$ the full spectrum is computed densely.  Eigenvalues keep their
sign and complex phase, since replacing $\lambda$ by $|\operatorname{Re}\lambda|$
would change Kemeny's sum; eigenvalues not computed by ARPACK are approximated
by zero, each contributing $1$ to $K$.

\subsection{Numerical Results}

\subsubsection{Metro: criticality is concentrated in three interchanges}
\label{sec:finding1}

Table~\ref{tab:resilience} ranks the metro stations.  Three stations stand
clearly apart: De\'ak Ferenc t\'er ($C_i=1.000$), where M1, M2 and M3 meet,
K\'alvin t\'er (0.619; M3/M4) and Keleti p\'alyaudvar (0.439; M2/M4).  They
carry the largest stationary departure flows (123, 60 and 43\,pax\,s$^{-1}$,
against about 15\,pax\,s$^{-1}$ at an ordinary station), and removing De\'ak
Ferenc t\'er alone reduces network efficiency by 49\%, K\'alvin t\'er by 32\%.
Below them the scores drop to 0.22--0.24 and are nearly flat; the order among
these stations is decided by the efficiency loss, which is largest for
stations close to the city centre (Arany J\'anos utca,
Corvin-negyed, Nyugati p\'alyaudvar).

\begin{table}[htbp]
\centering
\small
\begin{tabular}{rlrrrr}
\toprule
Rk & Station & $I_i$ (pax\,s$^{-1}$) & $V^{\mathrm{Kem}}$ & $V^{\mathrm{eff}}$ & $C_i$ \\
\midrule
1 & De\'ak Ferenc t\'er & 123.0 & 1.00 & 0.487 & 1.000 \\
2 & K\'alvin t\'er & 59.5 & 1.00 & 0.324 & 0.619 \\
3 & Keleti p\'alyaudvar & 43.3 & 1.00 & 0.096 & 0.439 \\
4 & Arany J\'anos utca & 15.0 & 1.00 & 0.181 & 0.237 \\
5 & Corvin-negyed & 14.9 & 1.00 & 0.158 & 0.232 \\
6 & Nyugati p\'alyaudvar & 15.0 & 1.00 & 0.148 & 0.231 \\
7 & F\H{o}v\'am t\'er & 15.2 & 1.00 & 0.131 & 0.230 \\
8 & Lehel t\'er & 15.1 & 1.00 & 0.122 & 0.227 \\
9 & Semmelweis Klinik\'ak & 14.7 & 1.00 & 0.128 & 0.224 \\
10 & Szent Gell\'ert t\'er - M\H{u}egyetem & 15.1 & 1.00 & 0.098 & 0.224 \\
11 & D\'ozsa Gy\"orgy \'ut & 15.1 & 1.00 & 0.098 & 0.224 \\
12 & G\"oncz \'Arp\'ad v\'arosk\"ozpont & 15.3 & 1.00 & 0.076 & 0.221 \\
\bottomrule
\end{tabular}
\caption{Metro: top-12 stations by criticality $C_i$ (all 48 stations evaluated).}
\label{tab:resilience}
\par\smallskip{\footnotesize $V^{\mathrm{Kem}}=1$: removal disconnects the metro chain.}
\forras{own computation from the BKK GTFS Schedule feed (CC0-1.0)}
\end{table}


The Kemeny component separates two groups.  For 36 of the 48 stations the
reduced chain is no longer irreducible: removing an intermediate station cuts
its line into pieces that cannot reach each other, and the fragmentation
penalty $V^{\mathrm{Kem}}=1$ applies.  The remaining 12 stations fall into two
groups.  Seven are termini (for example Kelenf\"old vas\'ut\'allom\'as,
\"Ors vez\'er tere, \'Ujpest-k\"ozpont), whose removal shortens a line and
does not increase $K$.  The other five (Ferenciek tere, Astoria, Blaha Lujza
t\'er, R\'ak\'oczi t\'er, II.~J\'anos P\'al p\'apa t\'er) lie on the inner
cycle De\'ak Ferenc t\'er -- K\'alvin t\'er -- Keleti p\'alyaudvar formed by
sections of M2, M3 and M4; removing one of them leaves an alternative route
round the cycle, and $K$ grows by only 6--26\%.  Kemeny's constant thus acts
mainly as a fragmentation indicator, which is what one expects of a network
of four long lines with three transfer stations: only the inner cycle provides
redundancy.

This is a structural statement about the timetable, not about observed
crowding: no fare-collection data enter the model.  The scheduled metro network
routes a disproportionate share of its stochastic passenger flow through three
interchange nodes, and this concentration is visible from the published GTFS
feed alone.

\subsubsection{Tram and bus}
\label{sec:finding2}

For the surface modes only the efficiency component is defined
(Section~\ref{sec:wellposed}), and removing a single stop changes the
efficiency of these large, sparse graphs very little: among the candidates the
largest reduction is 3.1\% for tram and 3.7\% for bus stops, and at the
high-throughput stops it is below 0.3\% (Table~\ref{tab:resilience-surface}).
The tram ranking is led by Sz\'ell K\'alm\'an t\'er M ($C_i=0.281$), which
combines by far the largest tram throughput (125\,pax\,s$^{-1}$) with a
measurable efficiency loss.  The following places are taken by stops of the
northern Pest tram corridor (\'Ujpest, Angyalf\"old), which have moderate
throughput but the largest efficiency losses (2.8--3.1\%).  Since the
efficiency of the tram and bus graphs is estimated from 200 source nodes,
differences of this size should be read with caution.  On the bus network the criticality scores are low
throughout ($C_i\le0.05$); the list is headed by bus stops at metro stations
(M\'oricz Zsigmond k\"ort\'er, K\H{o}b\'anya-Kispest) and at large bus
terminals (Buda\"orsi lak\'otelep, \'Ujpalota), i.e.\ by the feeder role of
buses to the metro and to the outer districts.  The small
efficiency losses reflect the tree-like, largely one-directional structure of
the surface graphs, in which a removed stop affects only paths through its own
branch; a joint multi-modal chain with walking transfers would be needed to
measure detours through other modes.

\begin{table}[htbp]
\centering
\small
\begin{tabular}{rllrrr}
\toprule
Rk & Stop & Mode & $I_i$ (pax\,s$^{-1}$) & $V^{\mathrm{eff}}$ & $C_i$ \\
\midrule
1 & Sz\'ell K\'alm\'an t\'er M & Tram & 124.7 & 0.0025 & 0.281 \\
2 & R\'akospalota-\'Ujpest vas\'ut\'allom\'as & Tram & 1.9 & 0.0313 & 0.118 \\
3 & Lehel t\'er M & Tram & 3.9 & 0.0143 & 0.117 \\
4 & Gy\"ongy\"osi utca & Tram & 2.1 & 0.0283 & 0.116 \\
5 & T\'el utca / Pozsonyi utca & Tram & 2.1 & 0.0279 & 0.115 \\
6 & Angyalf\"old vas\'ut\'allom\'as & Tram & 2.0 & 0.0282 & 0.115 \\
\midrule
1 & M\'oricz Zsigmond k\"ort\'er M & Bus & 142.7 & 0.0002 & 0.051 \\
2 & K\H{o}b\'anya-Kispest M & Bus & 300.4 & 0.0001 & 0.045 \\
3 & Buda\"orsi lak\'otelep & Bus & 82.6 & 0.0001 & 0.033 \\
4 & \'Ujpalota, Ny\'{\i}rpalota \'ut & Bus & 253.1 & 0.0000 & 0.031 \\
5 & K\H{o}b\'anya-Kispest M & Bus & 189.9 & 0.0001 & 0.030 \\
6 & \'Obuda, Bogd\'ani \'ut & Bus & 86.3 & 0.0001 & 0.026 \\
\bottomrule
\end{tabular}
\caption{Tram and bus: top-6 stops by criticality $C_i$ among the 100 candidates with the largest stationary throughput.}
\label{tab:resilience-surface}
\par\smallskip{\footnotesize The tram and bus chains are reducible, so $C_i$ uses the efficiency component only.}
\forras{own computation from the BKK GTFS Schedule feed (CC0-1.0)}
\end{table}


% =============================================================================
\section{Robustness-Driven Network Augmentation}
\label{sec:augment}

Section~\ref{sec:resilience} identifies the stations whose loss hurts the
network most.\footnote{[AUTHOR TO CONFIRM] In this section Claude (Anthropic)
was used beyond the level stated in the AI declaration: for drafting text,
the formulation of the propositions, the implementation in
\texttt{scripts/paper\_augment.py} and the literature search on planned
projects. The author has checked the derivations and the results.} The natural next question for a planner is the inverse one:
\emph{which new links would reduce that damage the most?}  This section turns
the node-removal analysis into a link-augmentation problem for the metro and
tram layers, solves it with an exact and fast shortest-path update, checks the
greedy solution against exhaustive search, and compares the suggested links
with the projects that are currently planned for Budapest.

\subsection{Problem Statement}

\textbf{Planning layers.}  For each mode $v\in\{\text{metro},\text{tram}\}$ we
work on a station-level directed graph $G=(S,E)$ with arc weights equal to the
generalised costs $C_{ij}$ of Eq.~\eqref{eq:cost}.  The metro layer is the
48-station graph of Section~\ref{sec:platforms}.  For the tram layer the
directional platforms are merged into \emph{stop areas}: platforms whose names
agree after removing parenthesised qualifiers and which lie within 400\,m of
each other form one node.  This turns the 651 tram platforms into 309 stop
areas.  The platform-level tram graph has 97 strongly connected components
(Section~\ref{sec:wellposed}), but the stop-area graph is strongly connected:
the reducibility of the tram chain is an artefact of one-directional platforms,
not of the network.

\textbf{Demand weights.}  Origin--destination pairs are weighted by the
gravity form of the E1 prior of Section~\ref{sec:prior},
\begin{equation}
  \omega_{ij}=\frac{p_ip_j}{\sum_{k\ne l}p_kp_l},\qquad
  p_i=\frac{\hat N_i(0)}{\sum_k\hat N_k(0)},\qquad i\ne j,
  \label{eq:omega}
\end{equation}
so that pairs of busy stations count more than pairs of lightly used ones.

\textbf{Candidate links.}  A candidate is an unordered pair $\{a,b\}$ of
stations that are not adjacent in $G$.  Adding it inserts the two arcs
$a\to b$ and $b\to a$ with cost
\begin{equation}
  c_{ab}=\hat c^{\,v}\,\ell_{ab},\qquad
  \hat c^{\,v}=\operatorname{median}_{(i,j)\in E}\frac{C_{ij}}{\ell_{ij}},
  \label{eq:newcost}
\end{equation}
where $\ell_{ab}$ is the great-circle distance and $\hat c^{\,v}$ the median
cost per metre of the existing links of the same mode ($0.060$\,s\,m$^{-1}$
for metro and $0.147$\,s\,m$^{-1}$ for tram).  The candidate set
$\mathcal C^v$ contains the pairs with $150\,\text{m}\le\ell_{ab}\le\ell^v_{\max}$
whose current cost is at least $1.5$ times the direct cost,
$\min(D_{ab},D_{ba})\ge1.5\,c_{ab}$, so that the link is a real shortcut.  We use
$\ell^{\mathrm{metro}}_{\max}=3$\,km and $\ell^{\mathrm{tram}}_{\max}=1.5$\,km.
A new tram link may not cross the Danube, since that would need a new bridge.
Metro links may cross it, as M2 and M4 already do.  This gives
$|\mathcal C^{\mathrm{metro}}|=229$ and $|\mathcal C^{\mathrm{tram}}|=579$.

\textbf{Vulnerability of an augmented network.}  Let $G_A$ denote $G$ with the
links of $A\subseteq\mathcal C^v$ added, $D^A$ its all-pairs cost matrix, and
$G_A-k$ the same network after the failure of station $k$.  With the
demand-weighted efficiency restricted to the OD pairs that do not involve $k$,
\begin{equation}
  E_k(H)=\sum_{\substack{i\ne j\\ i,j\ne k}}\frac{\omega_{ij}}{D^{H}_{ij}},
  \qquad 1/\infty:=0,
\end{equation}
the vulnerability of the augmented network to the loss of $k$ is
\begin{equation}
  V_k(A)=1-\frac{E_k(G_A-k)}{E_k(G_A)}\in[0,1].
  \label{eq:vk}
\end{equation}
$V_k(A)$ is the share of the (surviving) passengers' accessibility that is lost
when $k$ fails.  The demand that starts or ends at $k$ itself is excluded,
because no new link can serve it.  For $A=\emptyset$ and uniform $\omega$ this
is closely related to the efficiency component $V^{\mathrm{eff}}$ of
Section~\ref{sec:resilience}, but uses costs instead of hop counts and
excludes the pairs that start or end at $k$.

\textbf{Objective.}  Following the choice of an average and a worst-case
criterion, the robustness functional over a failure set $\mathcal F\subseteq S$
is
\begin{equation}
  J(A)=\alpha\,\bar V(A)+(1-\alpha)\,V_{\max}(A),\qquad
  \bar V(A)=\frac{1}{|\mathcal F|}\sum_{k\in\mathcal F}V_k(A),\quad
  V_{\max}(A)=\max_{k\in\mathcal F}V_k(A),
  \label{eq:J}
\end{equation}
with $\alpha=\tfrac12$.  The augmentation problem with a budget of $B$ links is
\begin{equation}
  A^\star_B\in\operatorname*{arg\,min}_{A\subseteq\mathcal C^v,\ |A|\le B} J(A).
  \label{eq:augprob}
\end{equation}
For the metro $\mathcal F$ contains all 48 stations.  For the tram it contains
the 100 stop areas with the largest prior demand, as in the node-removal sweep
of Section~\ref{sec:resilience}.

\subsection{Structural Results}

The following statements make the problem tractable and explain the results.

\begin{proposition}[Exact shortcut update]
\label{prop:shortcut}
Let $D$ be the all-pairs shortest-path cost matrix of a directed graph with
non-negative arc costs, and add the arcs $a\to b$ and $b\to a$ with cost
$c\ge0$.  The new cost matrix is
\begin{equation}
  D'_{ij}=\min\bigl\{D_{ij},\;D_{ia}+c+D_{bj},\;D_{ib}+c+D_{aj}\bigr\}.
  \label{eq:shortcut}
\end{equation}
The same formula applied to $D^{(-k)}$, the cost matrix of $G-k$, gives the cost
matrix of $G_{\{ab\}}-k$ for every $k\notin\{a,b\}$; for $k\in\{a,b\}$ the matrix
is unchanged.
\end{proposition}

\begin{proof}
Every expression on the right is the cost of a walk from $i$ to $j$ in the new
graph, so $D'_{ij}$ is at most the minimum.  Conversely, take a shortest
$i$--$j$ path in the new graph.  With non-negative costs we may assume that it
is simple, so it uses at most one of the two new arcs: using $a\to b$ and later
$b\to a$ (or the reverse) would visit $a$ twice.  If it uses no new arc its cost
is at least $D_{ij}$.  If it uses $a\to b$, its parts before and after this arc
are paths of the old graph from $i$ to $a$ and from $b$ to $j$, so its cost is
at least $D_{ia}+c+D_{bj}$; the case $b\to a$ is symmetric.  When $k$ has been
removed, a link with an endpoint at $k$ does not exist, and otherwise the same
argument holds in $G-k$.
\end{proof}

Applying \eqref{eq:shortcut} link by link gives the cost matrices for any set
$A$.  After a one-time computation of $D$ and of $D^{(-k)}$ for all
$k\in\mathcal F$ ($|\mathcal F|+1$ runs of Dijkstra's algorithm), evaluating a
candidate costs $O(|\mathcal F|\,n^2)$ elementary operations and requires no
further shortest-path computation.  In the implementation the update agrees
with a full recomputation to $10^{-13}$\,s.

\begin{proposition}[Extensions do not add redundancy]
\label{prop:pendant}
Let the augmentation attach a new subgraph $H$ (for example a line extension
with new stations) to $G$ at a single station $s$.  Then $D^A_{ij}=D_{ij}$ for all
$i,j\in S$ and, with the demand weights of the existing stations,
$V_k(A)=V_k(\emptyset)$ for all $k\in S$.
\end{proposition}

\begin{proof}
A path between $i,j\in S$ that enters $H$ has to enter and leave it through
$s$, so it contains a closed walk through $s$.  Removing that walk leaves a path
in $G$ that costs no more.  Hence all costs within $S$, with or without the
failure of any $k\in S$, are unchanged, and so are the efficiencies in
\eqref{eq:vk}.
\end{proof}

The planned M3 extension to Rákospalota-Újpest \citep{telex_m3_2026}, the
through-running of M2 onto the Gödöllő HÉV \citep{gov_h8h9} and a possible
M4 extension to Bosnyák tér are extensions of this kind.  They improve
coverage, but by Proposition~\ref{prop:pendant} they do nothing for the
robustness of the existing network.  If the new stations carry demand, the
failure of the attachment station $s$ also cuts them off, so $s$ becomes
\emph{more} critical.

\begin{proposition}[Hub-reinforcement trap]
\label{prop:trap}
If every link in $A$ has an endpoint at $k$, then $V_k(A)\ge V_k(\emptyset)$.
\end{proposition}

\begin{proof}
The links of $A$ disappear together with $k$, so $G_A-k=G-k$ and the
numerator of \eqref{eq:vk} is unchanged, while adding arcs can only shorten
paths, so $E_k(G_A)\ge E_k(G)$.
\end{proof}

New links into an already critical interchange raise the accessibility that is
lost when that interchange fails, so they make it more critical.  Robustness
improves through links that \emph{bypass} a hub, not through links that feed it.

\begin{remark}[Hardness and the greedy heuristic]
The closely related problem of adding $B$ shortcut edges to minimise the
average shortest-path distance is NP-hard \citep{meyerson2009};
\eqref{eq:augprob} adds failures and a worst-case term, and no exact
polynomial algorithm is known for it.  The reduction $J(\emptyset)-J(A)$ is
neither monotone (by Proposition~\ref{prop:trap} a link can increase some
$V_k$) nor submodular in general, since $J$ contains a maximum.  The
$(1-1/e)$ guarantee of the greedy algorithm for monotone submodular
maximisation \citep{nemhauser1978} therefore does not apply.  We use the greedy algorithm
and measure its gap empirically against exhaustive search for $B=2$.
\end{remark}

\begin{remark}[Periodicity]
\label{rem:periodic}
The metro chain is periodic because its support graph is bipartite
(Section~\ref{sec:wellposed}).  A new link between two stations of the same
colour class creates an odd cycle.  Since every metro link is used in both
directions, the chain also has cycles of length~2, so its period becomes
$\gcd(2,\text{odd})=1$.  Such a link therefore makes the chain aperiodic and
the spectral gap $g(P)$ positive, which makes the component $V^{\mathrm{gap}}$
of Section~\ref{sec:resilience} well defined again.
\end{remark}

\subsection{Algorithm}

\begin{figure}[H]
\centering
\fbox{\begin{minipage}{0.92\textwidth}
\textbf{Algorithm 2: Greedy robustness-driven augmentation}\\[4pt]
\textit{Require:} layer $G$, costs $C$, weights $\omega$, failure set
$\mathcal F$, candidates $\mathcal C$, budget $B$, committed links $A_0$
(empty, or a planned project).\\
\textit{Ensure:} links $a_1,\dots,a_B$ and $J$ after each step.\\[4pt]
\begin{enumerate}
  \item Compute $D$ and $D^{(-k)}$, $k\in\mathcal F$, by Dijkstra; apply
    \eqref{eq:shortcut} for the links of $A_0$.  $A\leftarrow A_0$.
  \item \textbf{for} $b=1,\dots,B$ \textbf{do}
  \begin{enumerate}
    \item \textbf{for each} $e\in\mathcal C\setminus A$: update $D$ and all
      $D^{(-k)}$ with $e$ by \eqref{eq:shortcut}, and evaluate $J(A\cup\{e\})$
      by \eqref{eq:vk}--\eqref{eq:J}.
    \item $a_b\leftarrow\arg\min_e J(A\cup\{e\})$;\quad $A\leftarrow A\cup\{a_b\}$.
  \end{enumerate}
  \item \textbf{return} $a_1,\dots,a_B$
\end{enumerate}
\end{minipage}}
\caption{Greedy robustness-driven link augmentation.}
\label{alg:augment}
\forras{own computation, \texttt{scripts/paper\_augment.py}}
\end{figure}

The cost is $O(B\,|\mathcal C|\,|\mathcal F|\,n^2)$: about
$1.3\times10^8$ operations for the metro with $B=5$, and about
$2.7\times10^{10}$ for the tram.  For the tram the candidate list is therefore
first reduced to the 100 best single links by $J(\{e\})$.  The metro run,
including the exhaustive check of all $\binom{229}{2}=26\,106$ pairs, takes
about two minutes on a laptop.

\subsection{Results: Metro}
\label{sec:aug-metro}

With the cost-based, demand-weighted definition the vulnerabilities of the
current metro reproduce the ranking of Section~\ref{sec:finding1}: the failure
of De\'ak Ferenc t\'er removes 44.5\% of the remaining accessibility and
K\'alvin t\'er 31.1\%, while no other station exceeds 16\%.  The worst-case
term of $J$ is thus set by the two central interchanges.

\begin{table}[htbp]
\centering
\footnotesize
\begin{tabular}{r>{\raggedright\arraybackslash}p{5.0cm}rrr>{\raggedright\arraybackslash}p{3.0cm}rr}
\toprule
$b$ & Added link & km & $\bar V$ & $V_{\max}$ & Worst failure & $J$ & $-\Delta J$ (\%) \\
\midrule
0 & \emph{current network} & --- & 0.070 & 0.445 & De\'ak Ferenc t\'er & 0.257 & --- \\
1 & R\'ak\'oczi t\'er -- Arany J\'anos utca & 1.73 & 0.067 & 0.328 & De\'ak Ferenc t\'er & 0.198 & 23.3 \\
2 & F\H{o}v\'am t\'er -- Opera & 1.77 & 0.060 & 0.208 & De\'ak Ferenc t\'er & 0.134 & 48.1 \\
3 & F\H{o}v\'am t\'er -- V\"or\"osmarty t\'er & 1.26 & 0.058 & 0.188 & K\'alvin t\'er & 0.123 & 52.1 \\
4 & H\H{o}s\"ok tere -- D\'ozsa Gy\"orgy \'ut & 1.53 & 0.046 & 0.185 & K\'alvin t\'er & 0.115 & 55.1 \\
5 & II. J\'anos P\'al p\'apa t\'er -- N\'epliget & 2.73 & 0.038 & 0.170 & De\'ak Ferenc t\'er & 0.104 & 59.7 \\
\bottomrule
\end{tabular}
\caption{Metro: greedy robustness-driven link additions. Each row adds one bidirectional link to all previous ones. $\bar V$ and $V_{\max}$ are the mean and maximum demand-weighted vulnerability over single-station failures, $J=\tfrac12(\bar V+V_{\max})$, and the last column is the cumulative reduction of $J$.}
\label{tab:aug-metro}
\forras{own computation from the BKK GTFS Schedule feed (CC0-1.0), \texttt{scripts/paper\_augment.py}}
\end{table}


Table~\ref{tab:aug-metro} shows the greedy plan.  The first three links,
R\'ak\'oczi t\'er -- Arany J\'anos utca (1.7\,km), F\H{o}v\'am t\'er -- Opera
(1.8\,km) and F\H{o}v\'am t\'er -- V\"or\"osmarty t\'er (1.3\,km), form a ring of
connections around De\'ak Ferenc t\'er and K\'alvin t\'er.  Each of them joins
two lines whose only connection today runs through one of the two hubs (M3 and
M4 meet only at K\'alvin t\'er; M1 reaches M4 only through De\'ak Ferenc t\'er).  Together they cut the
worst-case loss from 44.5\% to 18.8\% and reduce $J$ by
52.1\%.  The fourth and fifth links (H\H{o}s\"ok tere -- D\'ozsa
Gy\"orgy \'ut on the M1/M3 axis and II.~J\'anos P\'al p\'apa t\'er -- N\'epliget
on the M3/M4 axis) mainly lower the average $\bar V$, because they give the
outer sections of the lines a second connection.  In line with
Proposition~\ref{prop:trap}, none of the chosen links ends at De\'ak Ferenc t\'er
or K\'alvin t\'er.

\textbf{Greedy against exhaustive search.}  For $B=2$ the exhaustive optimum
(F\H{o}v\'am t\'er -- Bajcsy-Zsilinszky \'ut and Arany J\'anos utca --
Bajcsy-Zsilinszky \'ut) gives $J=0.127$, the greedy pair $J=0.134$.  The
greedy solution is within 5\% of the optimum, and both optima build the
same kind of bypass between the M1, M3 and M4 stations around the two hubs.

\textbf{Sensitivity.}  The first three greedy links are identical for
$\ell_{\max}=2$\,km, 3\,km and 4\,km, and for the pure worst-case objective
($\alpha=0$).  They depend on the weighting, however.  With the pure average
($\alpha=1$) the greedy algorithm connects the outer sections of the lines
instead (G\"oncz \'Arp\'ad v\'arosk\"ozpont -- Sz\'echenyi f\"urd\H{o},
V\"or\"osmarty utca -- Semmelweis Klinik\'ak, \'Ujbuda-k\"ozpont --
V\"or\"osmarty t\'er), and with uniform demand weights it chooses K\'alvin t\'er --
Opera, F\H{o}v\'am t\'er -- D\'eli p\'alyaudvar and H\H{o}s\"ok tere -- D\'ozsa
Gy\"orgy \'ut.  The inner bypass ring is therefore the answer to the worst-case
question.  A planner who cares only about the average failure would build
different links.

\begin{table}[htbp]
\centering
\footnotesize
\begin{tabular}{lrrcr}
\toprule
Metro network & $K(P)$ & $g(P)$ & periodic & fragmenting stations \\
\midrule
current network & 302.2 & 0.0000 & yes & 36 \\
greedy, $b=1$ & 286.3 & 0.0103 & no & 36 \\
greedy, $b=2$ & 262.9 & 0.0105 & no & 35 \\
greedy, $b=3$ & 259.1 & 0.0148 & no & 35 \\
greedy, $b=4$ & 220.2 & 0.0150 & no & 27 \\
greedy, $b=5$ & 193.2 & 0.0152 & no & 23 \\
M5 inner section & 237.4 & 0.0086 & no & 32 \\
M5 inner section + 3 next links & 201.7 & 0.0204 & no & 25 \\
\bottomrule
\end{tabular}
\caption{Markov-chain diagnostics of the augmented metro chain ($\lambda=1/300$\,s$^{-1}$). The last column counts the stations whose removal leaves a chain that is not irreducible.}
\label{tab:aug-chain}
\forras{own computation from the BKK GTFS Schedule feed (CC0-1.0), \texttt{scripts/paper\_augment.py}}
\end{table}


\textbf{Effect on the Markov chain.}  Table~\ref{tab:aug-chain} evaluates the
augmented networks with the same Wilson generator as in
Section~\ref{sec:ctmc}.  The first greedy link already makes the metro chain
aperiodic (Remark~\ref{rem:periodic}), so its spectral gap becomes positive.
Kemeny's constant falls from 302.2 to 193.2 after five links, i.e.\
a random passenger reaches a random destination in fewer steps.  The number of
stations whose loss disconnects the chain falls from 36 to 23.

\subsection{Results: Tram}
\label{sec:aug-tram}

The tram stop-area graph is much less vulnerable than the metro: the largest
loss, at the Lehel utca / R\'obert K\'aroly k\"or\'ut stop area where line 1 crosses
lines 12/14, is 8.9\%, followed by Sz\'ell K\'alm\'an t\'er (8.6\%).  The surface
network already has many parallel routes, which matches the small efficiency
losses of Section~\ref{sec:finding2}.

\begin{table}[htbp]
\centering
\footnotesize
\begin{tabular}{r>{\raggedright\arraybackslash}p{5.0cm}rrr>{\raggedright\arraybackslash}p{3.0cm}rr}
\toprule
$b$ & Added link & km & $\bar V$ & $V_{\max}$ & Worst failure & $J$ & $-\Delta J$ (\%) \\
\midrule
0 & \emph{current network} & --- & 0.024 & 0.089 & Lehel utca / R\'obert K\'aroly k\"or\'ut & 0.057 & --- \\
1 & V\'ag\'any utca / R\'obert K\'aroly k\"or\'ut -- Hun utca & 0.47 & 0.024 & 0.086 & Sz\'ell K\'alm\'an t\'er M & 0.055 & 3.5 \\
2 & Batthy\'any t\'er M+H -- D\'eli p\'alyaudvar M & 1.25 & 0.023 & 0.080 & Sz\'ell K\'alm\'an t\'er M & 0.052 & 8.6 \\
3 & Oktogon M -- Lehel t\'er M & 1.30 & 0.021 & 0.080 & Sz\'ell K\'alm\'an t\'er M & 0.051 & 10.7 \\
4 & Blaha Lujza t\'er M -- K\'alvin t\'er M & 1.07 & 0.020 & 0.080 & Sz\'ell K\'alm\'an t\'er M & 0.050 & 12.3 \\
5 & Budafoki \'ut / Domb\'ov\'ari \'ut -- Budafoki \'ut / Szer\'emi sor & 0.47 & 0.019 & 0.079 & Sz\'ell K\'alm\'an t\'er M & 0.049 & 13.1 \\
\bottomrule
\end{tabular}
\caption{Tram: greedy robustness-driven link additions. Each row adds one bidirectional link to all previous ones. $\bar V$ and $V_{\max}$ are the mean and maximum demand-weighted vulnerability over single-station failures, $J=\tfrac12(\bar V+V_{\max})$, and the last column is the cumulative reduction of $J$.}
\label{tab:aug-tram}
\forras{own computation from the BKK GTFS Schedule feed (CC0-1.0), \texttt{scripts/paper\_augment.py}}
\end{table}


The greedy tram plan (Table~\ref{tab:aug-tram}) consists of short gap
closures.  (1)~A 0.5\,km link V\'ag\'any utca / R\'obert K\'aroly k\"or\'ut --
Hun utca joins line 1 to lines 12/14 away from their only crossing and removes
the worst case.  (2)~Batthy\'any t\'er -- D\'eli p\'alyaudvar (1.2\,km) connects the
riverside lines 19/41 with lines 17, 56, 59 and 61 without passing through
Sz\'ell K\'alm\'an t\'er.  (3)~Oktogon -- Lehel t\'er (1.3\,km)
joins the Grand Boulevard lines 4/6 to lines 12/14.  (4)~Blaha Lujza t\'er --
K\'alvin t\'er (1.1\,km) joins lines 4/6 and 28/37 to lines 47/49 along the
inner boulevard.  (5)~A 0.5\,km gap on Budafoki \'ut joins line 1 to line 4.
It lies on the corridor of the Budafoki \'ut tram line, which is currently
being designed \citep{bkk_budafoki2026}.  Five links reduce
$J$ by 13.1\%.  The worst case moves to Sz\'ell K\'alm\'an t\'er,
whose loss the five links reduce only from 8.6\% to 7.9\%.

\subsection{Planned Projects and Suggested Next Links}
\label{sec:aug-planned}

\begin{table}[htbp]
\centering
\footnotesize
\begin{tabular}{l>{\raggedright\arraybackslash}p{4.6cm}rrrrrr}
\toprule
Layer & Project & Links & km & $\bar V$ & $V_{\max}$ & $-\Delta J$ (\%) & greedy, same $b$ (\%) \\
\midrule
Metro & M5 inner section (K\'alvin t\'er--Astoria--Oktogon--Lehel t\'er) & 3 & 3.2 & 0.052 & 0.320 & $27.8$ & 52.1 \\
Tram & Pesti fon\'od\'o (De\'ak t\'er--Nyugati t\'er--Lehel t\'er) & 2 & 2.3 & 0.021 & 0.096 & $-2.7$ & 8.6 \\
Tram & Th\"ok\"oly \'ut tram (Keleti pu.--Bosny\'ak t\'er), discussed & 1 & 3.2 & 0.023 & 0.089 & $1.1$ & 3.5 \\
\bottomrule
\end{tabular}
\caption{Planned and discussed projects evaluated on the same robustness functional. The last column is the reduction achieved by the greedy plan with the same number of links.}
\label{tab:aug-planned}
\par\smallskip{\footnotesize Pure extensions (M3 to Rákospalota-Újpest, M2--Gödöllő HÉV, M4 to Bosnyák tér) are not listed: by Proposition~\ref{prop:pendant} they leave every $V_k$ unchanged.}
\forras{own computation from the BKK GTFS Schedule feed (CC0-1.0), \texttt{scripts/paper\_augment.py}}
\end{table}


Table~\ref{tab:aug-planned} evaluates the projects that are currently planned
or discussed with the same functional.

\textbf{M5 inner section.}  The inner section of the north--south regional rail
(M5) would run from Boráros tér through K\'alvin t\'er, Astoria and Oktogon to
Lehel t\'er \citep{wiki_m5}.  Represented by its three links between existing
metro stations (the new station at Klauz\'al t\'er is not a node of the current
graph), it lowers $J$ by 27.8\%.  This is roughly half of what the
greedy plan achieves with the same number of links (52.1\%).  The
reason is Proposition~\ref{prop:trap} at work: the M5 passes through K\'alvin
t\'er, the second most critical station, and makes it the new worst case
(32.0\%).  The M5 is still the best of all evaluated plans for the intact
network: it raises the demand-weighted efficiency by
5.8\%, against 1.3\% for the greedy links.  M5 is
therefore primarily a capacity and accessibility project, not a redundancy project.

\textbf{Pesti fonód\'o.}  The new tram connection Deák tér -- Nyugati tér --
Lehel tér joins lines 47/49 to lines 12/14 \citep{telex_fonodo2024,bkk_fonodo2025}.
It lowers the average vulnerability $\bar V$ by
13.1\%, against 2.1\% for the first two greedy links together, and raises the
demand-weighted efficiency by 3.5\%.  At the same time it
routes additional paths through the Lehel utca / R\'obert K\'aroly k\"or\'ut
stop area, whose failure becomes more damaging (from 8.9\% to 9.6\%), so $J$
as a whole does not improve.  The Th\"ok\"oly \'ut tram (Keleti -- Bosny\'ak
t\'er), which has been discussed but not planned, has only a small effect.

\begin{table}[htbp]
\centering
\footnotesize
\begin{tabular}{l>{\raggedright\arraybackslash}p{5.4cm}rrr>{\raggedright\arraybackslash}p{3.0cm}r}
\toprule
Layer & Project / next link & km & $\bar V$ & $V_{\max}$ & Worst failure & $-\Delta J$ (\%) \\
\midrule
Metro & \emph{M5 inner section only} & --- & 0.052 & 0.320 & K\'alvin t\'er & $27.8$ \\
 & +1: Astoria -- Corvin-negyed & 1.21 & 0.050 & 0.181 & K\'alvin t\'er & $55.1$ \\
 & +2: Szent Gell\'ert t\'er - M\H{u}egyetem -- D\'eli p\'alyaudvar & 2.92 & 0.043 & 0.139 & Corvin-negyed & $64.7$ \\
 & +3: V\"or\"osmarty utca -- Semmelweis Klinik\'ak & 2.88 & 0.039 & 0.112 & Lehel t\'er & $70.6$ \\
\midrule
Tram & \emph{Pesti fon\'od\'o only} & --- & 0.021 & 0.096 & Lehel utca / R\'obert K\'aroly k\"or\'ut & $-2.7$ \\
 & +1: V\'ag\'any utca / R\'obert K\'aroly k\"or\'ut -- Hun utca & 0.47 & 0.021 & 0.087 & Sz\'ell K\'alm\'an t\'er M & $5.3$ \\
 & +2: Blaha Lujza t\'er M -- K\'alvin t\'er M & 1.07 & 0.020 & 0.086 & Sz\'ell K\'alm\'an t\'er M & $6.7$ \\
 & +3: Budafoki \'ut / Domb\'ov\'ari \'ut -- Budafoki \'ut / Szer\'emi sor & 0.47 & 0.019 & 0.085 & Sz\'ell K\'alm\'an t\'er M & $7.5$ \\
\bottomrule
\end{tabular}
\caption{Suggested next links once the planned project is built: greedy additions conditional on the project's links. $-\Delta J$ is measured against the current network.}
\label{tab:aug-next}
\forras{own computation from the BKK GTFS Schedule feed (CC0-1.0), \texttt{scripts/paper\_augment.py}}
\end{table}


\textbf{What to add next.}  Table~\ref{tab:aug-next} answers the planning
question directly: if the planned project is built, which links should follow?
The rows are the output of Algorithm~\ref{alg:augment} with the project's
links committed ($A_0\ne\emptyset$).  For the metro, the M5 is best complemented by Astoria -- Corvin-negyed
(1.2\,km), which joins M2 and M5 to M3 without passing through K\'alvin t\'er
and lowers the worst case from 32.0\% to 18.1\%.  It is followed by two longer
outer links, Szent Gell\'ert t\'er -- D\'eli p\'alyaudvar (M4 to M2 on the Buda
side, 2.9\,km) and V\"or\"osmarty utca -- Semmelweis Klinik\'ak (M1 to M3,
2.9\,km).  M5 with these three links reduces $J$ by 70.6\%, more than five
greedy links on the current network (59.7\%), and lowers Kemeny's constant to
201.7 (Table~\ref{tab:aug-chain}).  For the tram, the Pesti fonód\'o is best
complemented by the short 1 -- 12/14 link at V\'ag\'any utca -- Hun utca, which
removes the new weak point at Lehel utca / R\'obert K\'aroly k\"or\'ut, and then
by Blaha Lujza t\'er -- K\'alvin t\'er and the Budafoki \'ut gap.  With these
links the fonód\'o turns from a small loss ($-2.7\%$) into a gain of 7.5\% in
$J$, while keeping its large gain in average accessibility.

\subsection{Interpretation and Limits}

Three conclusions follow from the analysis.  First, the Budapest metro is
fragile at exactly two points, and the most effective remedy is a small set
of bypass links that join the radial lines \emph{without} passing through
De\'ak Ferenc t\'er or K\'alvin t\'er.  Three links of about 1.3--1.8\,km
halve $J$.  Second, extensions of existing lines, which account for most of
the current metro plans, do not make the existing network more robust
(Proposition~\ref{prop:pendant}).  Third, planned projects that serve an
existing hub (M5 at K\'alvin t\'er, the Pesti fonód\'o at Lehel t\'er) improve
average accessibility but shift the worst case to that hub.  Each of them
should be combined with the complementary links of Table~\ref{tab:aug-next}.

The analysis is structural.  The costs of new links are the median costs per
metre of the mode, so construction feasibility, alignment and actual cost
are not modelled, and the suggested links should be read as corridors
(``connect these two lines between these two areas''), not as alignments.  The
station graph has no transfer penalty, so it overstates the value of links
that create paths with an extra change of line.  The demand weights come from
the E1 prior, and the sensitivity analysis shows that the average-oriented
recommendations depend on them.  Finally, metro and tram are treated as
separate layers.  In a joint layer with walking transfers, a tram line can
replace a failed metro station, so the joint vulnerabilities, and possibly the
recommended links, would be different.

% =============================================================================
\section{Limitations and Future Work}
\label{sec:limitations}

We enumerate the principal limitations of the current model.

\textbf{Idealised passenger behaviour.}  Independence assumption, single-leg
transitions, and origin-constrained (not doubly-constrained) gravity model.

\textbf{Time-varying supply.}  The time-homogeneous CTMC is valid only within
a narrow service window where $f^v_i$ is approximately constant.

\textbf{Capacity constraints.}  The uncapacitated formulation ignores finite
vehicle capacity.  A capacity-dampened rate
$\mu^v_i(N^v_i)=\mu^v_i\cdot\max(0,1-N^v_i/(K^v|\mathcal{N}^v_{\mathrm{in}}(i)|))$
is tractable for SSA but lacks a closed-form stationary distribution.

\textbf{Absence of validation data.}  No AFC or APC data were available.
Parameters $\lambda^v$ and $\gamma$ are assumed, not estimated, and the prior
uses the service-proxy estimator E1.  Absolute passenger counts are therefore
illustrative; the structural results (which stops dominate the ranking) are the
robust part of the analysis.

\textbf{Absorbing terminals.}  Stops without peak departures act as absorbing
states and accumulate 19\% of the passengers within 15 minutes
(Section~\ref{sec:results}).  An explicit trip-end (alighting) process would
remove this artefact and allow longer horizons.

\textbf{Mode coverage.}  Trolleybus and suburban rail are coded in the current
feed with \texttt{route\_type} values outside the framework's mode table and
are not part of the analysed network.

\textbf{Markov assumption.}  Real service headways follow scheduled
(deterministic) distributions, not exponential ones.  Phase-type or
semi-Markov extensions \citep{norris1997} can generalise this.

\textbf{Single-mode assignment.}  Mode switching mid-journey is not modelled.

\textbf{Per-mode resilience analysis.}  Node removal is evaluated within
each modal chain, so the ranking captures the role of a stop within its own
mode but not the loss of transfer opportunities between modes.  A joint
multi-modal generator $Q=\bigoplus_v Q^v$ with inter-modal transfer edges is
required for that.

\subsection{Future Work}

\begin{enumerate}
  \item \textbf{AFC calibration:} Estimate $\lambda^v$ from tap-in/tap-out
    distributions when farecard data are available.
  \item \textbf{GTFS-RT integration:} Update $f^v_i$ in real time from
    vehicle-position feeds.
  \item \textbf{Doubly-constrained formulation:} Implement IPF balancing
    for trip distribution.
  \item \textbf{Capacity-dampened SSA:} Nonlinear propensities with finite
    vehicle capacity.
  \item \textbf{Joint and costed augmentation:} Repeat the link-augmentation
    analysis of Section~\ref{sec:augment} on a joint metro--tram layer with
    walking transfers and transfer penalties, with construction costs per
    kilometre in place of a cardinality budget.
  \item \textbf{Regional comparison:} Apply to Vienna (WL), Prague (DPP),
    and Warsaw (ZTM) GTFS feeds.
  \item \textbf{Joint multi-modal criticality:} Node removal on the combined
    generator $Q=\bigoplus_v Q^v$ to recover cross-modal interchange
    significance.
  \item \textbf{Trip-end process and full mode coverage:} Add alighting as
    an explicit outflow and include trolleybus and suburban rail.
\end{enumerate}

% =============================================================================
\section{Summary and Conclusions}
\label{sec:conclusion}

We have developed a multi-modal, entropy-maximising Markov-chain
passenger-flow model for the BKK Budapest GTFS network and implemented it in
Python using sparse linear algebra throughout.  The model operates on three
complementary stochastic layers---a DTMC whose transition matrix is the Wilson
softmax, a CTMC whose generator is derived from the same entropy kernel with
GTFS-calibrated departure intensities, and an exact stochastic realisation via
the Gillespie SSA.

\textbf{Theoretical contributions.}  We derive the
CTMC generator from entropy-maximising flows with correct dimensional analysis
and sparse graph structure.  We show precisely how the DTMC embedded jump chain
relates to the CTMC, and characterise the conditions for ergodicity.

\textbf{Empirical findings.}  On the BKK feed (5\,177 stops, 6\,450 scheduled
links, four modes) the model conserves passengers to machine precision in the
KFE and exactly in the integer SSA and $\tau$-leap layers.  The mean field is an
accurate description at station level: exact SSA ensembles on the metro agree
with it within 0.8\% and fluctuate at the Poisson level $N^{-1/2}$, and
$\tau$-leaping on the full network agrees within 2.4\%.  Metro stations mix
within a few headways, while terminal stops without departures absorb a growing
share of passengers, which limits the time-homogeneous model to horizons of the
order of a trip.  In the resilience analysis the spectral gap turns out to be
uninformative on every modal chain (the metro chain is periodic, the surface
chains reducible), whereas Kemeny's constant works as a fragmentation
indicator.  Criticality in the metro is concentrated in the three interchanges
De\'ak Ferenc t\'er, K\'alvin t\'er and Keleti p\'alyaudvar, and only the inner
cycle formed by M2, M3 and M4 offers redundancy against the loss of a station.

\textbf{Network augmentation.}  Inverting the failure analysis gives concrete
planning suggestions.  The metro is fragile at two interchanges, and three links
that join the radial lines around De\'ak Ferenc t\'er and K\'alvin t\'er,
without passing through them, halve the robustness functional; greedy search
is within 5\% of the exhaustive optimum.  Line extensions provably add no
redundancy, and planned projects that run through a hub (M5 at K\'alvin t\'er,
the Pesti fonód\'o at Lehel t\'er) improve average accessibility but move the
worst case to that hub, so they should be complemented by the bypass links
identified in Section~\ref{sec:aug-planned}.

\textbf{Limitations.}  Without AFC or APC data, absolute passenger counts and
$\lambda^v$ parameters cannot be validated; the sensitivity sweep shows,
however, that the default $\lambda^v$ lies in a regime where the stationary
concentration is insensitive to it.  Capacity constraints, time-varying supply,
trip ends and multi-leg journeys are absent from the current formulation.

\textbf{Code and data availability.}  The full Python implementation is at
\url{https://github.com/momarci13/bkk-markov-passenger-flow} (MIT licence).  BKK GTFS:
\url{https://opendata.bkk.hu} (CC0-1.0).

% =============================================================================

\clearpage
\bibliographystyle{plainnat}
\bibliography{hivatkozasok}
\addcontentsline{toc}{section}{References}

\appendix
\renewcommand{\thesection}{\Alph{section}}
\renewcommand{\thesubsection}{\thesection.\arabic{subsection}}
\section*{Appendices}
\addcontentsline{toc}{section}{Appendices}
\setcounter{section}{0}
\section{Reproducibility}
\label{app:repro}

\subsection{Software Versions}

\begin{table}[htbp]
\centering
\small
\begin{tabular}{ll}
\toprule
Package & Version \\
\midrule
Python    & 3.11 \\
NumPy     & 2.4 \\
SciPy     & 1.17 \\
pandas    & 3.0 \\
NetworkX  & 3.6 \\
\bottomrule
\end{tabular}
\caption{Software environment used for the results in this paper.}
\forras{own computation, \texttt{bkk/} framework}
\end{table}

\subsection{Reproducibility Checklist}

\begin{enumerate}
  \item \texttt{pip install -e .}
  \item Download the feed from \url{https://opendata.bkk.hu} to
    \texttt{data/budapest\_gtfs.zip}.
  \item \texttt{python scripts/tdk\_rerun.py -{}-zip data/budapest\_gtfs.zip}.
    This runs the network construction, generators, KFE, $\tau$-leap and the
    node-removal sweep (\texttt{data/resilience\_scores.csv}), then
    \texttt{scripts/paper\_dynamics.py} (KFE trajectories, SSA and $\tau$-leap
    ensembles, $\lambda$ sweep; \texttt{data/paper\_dynamics.json}) and
    \texttt{scripts/paper\_tables.py}, which writes every table fragment in
    \texttt{paper/tables/}, and finally \texttt{scripts/paper\_augment.py}
    (Section~\ref{sec:augment}; \texttt{data/augment\_results.json} and the
    \texttt{t\_aug\_*} tables, about 35 minutes).  On Windows, \texttt{futtatas.ps1} performs all
    steps and builds the PDF.
\end{enumerate}
Random streams are seeded (seed 42 for the pipeline, fixed seeds per
realisation for the ensembles), so every number in
Sections~\ref{sec:results}--\ref{sec:resilience} is reproducible.


\subsection{Numerical Validation Tests}
\label{app:tests}

The package ships 91 \texttt{pytest} tests, all passing.  Among them:
\begin{itemize}
  \item Row sums of $Q^v$ equal zero to $<10^{-10}$ in $\ell^\infty$ norm.
  \item Row sums of $P^v$ equal one to $<10^{-14}$ in $\ell^\infty$ norm.
  \item $\|N^v(T)\|_1=\|N^v(0)\|_1$ (KFE conservation) to $<10^{-6}$ relative error.
  \item SSA: integer-valued state vector throughout; no negative entries.
  \item Prior: $\sum_i\hat{N}_i(0)=N_{\mathrm{peak}}$ to $<10^{-8}$ relative error.
  \item Metro platforms are aggregated to parent stations and the metro chain
    is irreducible with non-zero departure hazards.
  \item Eigenvalues keep sign and phase; reducible chains report an undefined
    Kemeny constant.
  \item The shortcut update \eqref{eq:shortcut} agrees with a full Dijkstra
    recomputation, with and without failures, on random graphs; pendant
    extensions leave every $V_k$ unchanged (Proposition~\ref{prop:pendant});
    links into $k$ never reduce $V_k$ (Proposition~\ref{prop:trap}).
\end{itemize}

\subsection*{Notation Summary}

\begin{table}[htbp]
\centering
\small
\begin{tabular}{lll}
\toprule
Symbol & Meaning & Unit / Range \\
\midrule
$n$          & Number of stops               & $5{,}177$ \\
$m$          & Number of modes               & 4 \\
$R$          & Total admissible edges        & $6{,}450$ \\
$C^v_{ij}$   & Generalised cost              & s \\
$\lambda^v$  & Entropy--cost multiplier      & s$^{-1}$ \\
$\pi^v_{j|i}$& Wilson destination kernel     & $[0,1]$, $\sum_j=1$ \\
$P^v_{ij}$   & DTMC transition               & $[0,1]$, row-stochastic \\
$q^v_{ij}$   & CTMC off-diagonal rate        & s$^{-1}$ \\
$\mu^v_i$    & Per-passenger departure hazard& s$^{-1}$ \\
$N^v_i(t)$   & Passenger count               & $\mathbb{Z}_{\ge0}$ \\
$\pi^v$      & Stationary distribution       & $[0,1]^n$, $\sum=1$ \\
$K(P)$       & Kemeny's constant             & steps \\
$g(P)$       & Spectral gap                  & $[0,1]$ \\
$E(P)$       & Network efficiency            & dimensionless \\
$I_i$        & Stationary departure exposure & pax\,s$^{-1}$ \\
$C_i$        & Stop criticality              & $[0,1]$ \\
$\omega_{ij}$& OD demand weight              & $[0,1]$, $\sum=1$ \\
$c_{ab}$     & Cost of a new link $a$--$b$   & s \\
$V_k(A)$     & Vulnerability to failure of $k$ after adding $A$ & $[0,1]$ \\
$J(A)$       & Robustness functional         & $[0,1]$ \\
\bottomrule
\end{tabular}
\caption{Principal symbols and units.}
\forras{own computation, \texttt{bkk/} framework}
\end{table}

% =============================================================================

% ----------------------------------------------------------------------
\section{Declaration on the Use of Artificial Intelligence}
% ----------------------------------------------------------------------
As required by the call for papers, I declare which artificial-intelligence
tools were used in preparing this paper, and for what purpose.

\begin{description}[leftmargin=0pt,style=nextline,nosep]
\item[Tool and version] Claude (Anthropic), Opus 5.5 model
\item[Literature search and topic mapping] Yes, for the related-work
overview in Section~\ref{sec:intro}.
\item[Language checking and proofreading] No.
\item[Code review and debugging] Yes, for the analysis code in
\texttt{bkk/} and \texttt{scripts/}.
\item[Substantive text generation] No. The research question, the mathematical
derivations, the implementation decisions and the interpretation of the results
are my own work.
\end{description}

\medskip
\noindent Where the depth of tool use in a given chapter or section differed
from the above, a footnote attached to the first sentence of that part says so.

% ----------------------------------------------------------------------
\section{Declaration of Originality}
% ----------------------------------------------------------------------
I declare that this paper is original: it is not identical to any paper of mine
submitted to another section, to any paper submitted to another TDK conference,
to a master's thesis submitted and defended, or to any published or forthcoming
publication.

\vspace{1.5cm}
\noindent\hrulefill\hspace{1cm}\hrulefill

\noindent date \hspace{5.6cm} signature of the author


\end{document}
